%================= 画像なし・最小構成（UTF-8） =================
\documentclass[a4paper,11pt,draft]{jsarticle}
% 余白と数式・表・URL
\usepackage[top=20mm,bottom=20mm,left=22mm,right=22mm]{geometry}
\usepackage{amsmath,amssymb}
\usepackage{booktabs,longtable,multirow}
\usepackage{url}

% 画像挿入用
\usepackage[dvipdfmx]{graphicx}
\usepackage[dvipdfmx]{xcolor}
\usepackage{subcaption}
\graphicspath{{figures/}}

% しおりの日本語対応（PDFしおりで文字化けする場合）
\usepackage[dvipdfmx]{hyperref}
\usepackage{pxjahyper}

\usepackage{tabularx}  % 全幅を使う表
\usepackage{booktabs}  % きれいな横罫線（\toprule など）
\usepackage{float}
\usepackage{cleveref}
\usepackage{acro}
\usepackage{graphicx}

\usepackage{numprint}
\npthousandsep{,} % 3桁区切り（点にしたいなら { . }）

\usepackage{hyperref}
\usepackage{xurl} % 長いURLの改行がうまくなる

\usepackage{listings}
\usepackage{jvlisting} % ← 日本語を含む listing 対策（重要）
\usepackage{tikz}

\usepackage{wrapfig}

% preamble に追加
\usepackage{seqsplit}
\usepackage{array}
\newcolumntype{P}[1]{>{\raggedright\arraybackslash}p{#1}}

\usepackage{titling}

% タイトル・著者・日付の見た目をここで指定
\pretitle{\begin{center}\huge\bfseries}
\posttitle{\par\end{center}\vskip 1.4em}

\preauthor{\begin{center}\Large}
\postauthor{\par\end{center}\vskip 0.9em}

\predate{\begin{center}\large}
\postdate{\par\end{center}\vskip 1.5em}



\lstset{
  basicstyle=\ttfamily\small,
  frame=single,
  breaklines=true,
  columns=fullflexible,
  showstringspaces=false
}





\newcommand{\repo}{\url{https://github.com/tomoki0715siu/RiNALMo}}
\newcommand{\repofile}[1]{\texttt{#1}（\repo{} ）}





\DeclareAcronym{CLS}{
  short = CLS ,
  long  = classification token
}
\DeclareAcronym{EOS}{
  short = EOS ,
  long  = end-of-sequence token
}
\DeclareAcronym{aaRS}{
  short = aaRS ,
  long  = aminoacyl-tRNA synthetase
}
\DeclareAcronym{ncRNA}{
  short = ncRNA ,
  long  = non-coding RNA
}
\DeclareAcronym{tRNA}{
  short = tRNA ,
  long  = transfer RNA
}
\DeclareAcronym{mRNA}{
  short = mRNA ,
  long  = messenger RNA
}
\DeclareAcronym{attnlrp}{
  short = AttnLRP ,
  long  = Attention-Aware Layer-wise Relevance Propagation
}
\DeclareAcronym{MHA}{
  short = MHA ,
  long  = Multi-Head Attention
}
\DeclareAcronym{FFN}{
  short = FFN ,
  long  = Feed-Forward Network
}
\DeclareAcronym{Rope}{
  short = Rope ,
  long  = Rotary Positional Embedding
}
\DeclareAcronym{MLM}{
  short = MLM ,
  long  = Masked Language Modeling
}
\DeclareAcronym{CCoS}{
  short = CCoS ,
  long  = centered cosine similarity
}
\DeclareAcronym{arc}{
  short = arc ,
  long  = Archaea
}
\DeclareAcronym{bac}{
  short = bac ,
  long  = Bacteria
}
\DeclareAcronym{euk}{
  short = euk ,
  long  = Eukaryota
}



\title{%
RNA言語モデルを用いた\\tRNA分類と内部表現解析\\[0.8em]
{\LARGE\mdseries Classification of tRNA\\ Using a Genomic Language Model and Analysis of Internal Representations}%
}

\author{北海道大学 情報エレクトロニクス学科 生体情報コース\\
情報生物学研究室\\
横川知樹}

\date{2026年3月}



\begin{document}
% 表紙のみ（1ページ目）
\maketitle
\thispagestyle{empty} % ページ番号非表示

\clearpage
\null
\thispagestyle{empty}
\clearpage % 要旨を新しいページから始めたいなら（不要なら消してOK）


\phantomsection
\tableofcontents
\clearpage

% jsarticle のデフォルト見出し名を好みに変更（任意だけどおすすめ）
\renewcommand{\listfigurename}{図一覧}
\renewcommand{\listtablename}{表一覧}

% 「目次に載せるだけ」用
\newcommand{\addtotoc}[2]{%
  \clearpage
  \phantomsection
  \addcontentsline{toc}{#1}{#2}%
}

% --- 図一覧 ---
\addtotoc{section}{\listfigurename}
\listoffigures

% --- 表一覧 ---
\addtotoc{section}{\listtablename}
\listoftables

% 略語一覧
\addtotoc{section}{略語一覧}
\printacronyms[name=略語一覧, heading=section*]

\clearpage
\pagestyle{plain}

\section{序論}

\subsection{研究背景}
\subsubsection{tRNAとは}

\ac{tRNA}は、遺伝情報の翻訳過程において中心的な役割を果たす\ac{ncRNA}である。
\ac{mRNA}上の遺伝暗号を読み取り、対応するアミノ酸をリボソームへ運搬することで、遺伝情報からタンパク質への変換を媒介する分子アダプターとして機能する。
\ac{tRNA}は一般に70〜90ヌクレオチドからなる比較的小さなRNA分子であり、二次構造としてクローバーリーフ型、三次構造としてL字型の立体構造(図\ref{fig:tRNAstructure})を形成する。
二次構造の特徴的な構造は、acceptor arm,D arm,anticodon arm,T arm,variable region(VR)といった保存された領域から構成されている。
特に、3'末端のCCA配列はアミノ酸が共有結合する部位であり、全てのtRNAに保存されている普遍的な特徴である。
pppp
\ac{tRNA}の機能において最も重要な要素の一つがアンチコドンである。アンチコドンはtRNA分子のanticodon arm上に位置する3つの連続したヌクレオチドであり、\ac{mRNA}上のコドンと相補的な塩基対を形成することで、遺伝暗号の解読を可能にする。

\begin{figure}[H]
  \centering
  \includegraphics[width=0.8\textwidth]{tRNA_structure.png}
  \caption{tRNAの二次構造(A)および三次構造(B)(\cite{giege2023trna} の Figure 1 を改変せず転載)}
  \label{fig:tRNAstructure}
\end{figure}



\subsubsection{tRNAのアミノ酸特異性}
\label{subsec:tRNA_aaRS}
各\ac{tRNA}への正確なアミノ酸の付加は、\ac{aaRS}によって触媒される。\ac{aaRS}は原則として20種類の標準アミノ酸それぞれに対応する酵素が存在し、対応する\ac{tRNA}を高度に識別して正しいアミノ酸を結合させる。

\ac{tRNA}のアミノ酸特異性を決定する要因は複雑かつ多面的である。アンチコドン配列は重要な識別要素の一つであるが、それのみでは十分ではなく、\ac{aaRS}は\ac{tRNA}分子全体の構造的特徴を認識する。
特に、acceptor arm,D arm,variable regionなどの領域に存在する特定のヌクレオチドが、\ac{aaRS}との相互作用において重要な決定因子として機能することが知られている\cite{giege2023trna}。
例えば、\ac{tRNA}$^{\mathrm{Ala}}$では acceptor arm に存在する G3:U70 塩基対が決定的な識別要素であり、アンチコドン配列を変化させてもこの塩基対が保存されていればアラニンが付加されることが示されている\cite{naganuma2014selective}。
これらの識別決定因子は、同一アミノ酸に対応する\ac{tRNA}であっても生物ドメイン間(\ac{arc}, \ac{bac}, \ac{euk})で異なる場合がある。
\subsubsection{tRNA研究の意義}
このように、\ac{tRNA}のアミノ酸特異性は多層的な要因が統合的に作用することで決定されている。配列情報からこれらの機能的特徴を理解することは、翻訳機構の本質的理解だけでなく、遺伝暗号の進化、tRNA遺伝子の同定、人工tRNAの設計といった応用的観点からも重要である。

\subsection{先行研究と課題}

\ac{tRNA}配列の解析において、従来は比較配列解析や構造予測が主要な手法であった。
tRNAscan-SE\cite{chan2021trnascan}などのツールは、ゲノム配列からtRNA遺伝子を同定するために広く用いられており、保存された配列モチーフや二次構造の特徴に基づいて\ac{tRNA}を検出する。
これらの手法は高い感度と特異度を持つが、既知の構造的特徴に依存するため、非典型的な\ac{tRNA}や新規の\ac{tRNA}配列の検出には限界がある。また、GtRNAdb\cite{chan2016gtrnadb}やtRNAdb-CE\cite{abe2014trnadb}といったデータベースは、実験的に検証された\ac{tRNA}配列を集約し、アノテーション情報を提供しているが、配列とアミノ酸特異性の関係性を体系的に解釈する枠組みは提供していない。

実験的な変異導入やin vitro選択実験により、各アミノ酸に対応する\ac{tRNA}の識別要素が個別に明らかにされてきた。しかし、これらの研究は労力と時間を要し、全ての\ac{tRNA}種について網羅的に解析することは現実的でない。また、複数の塩基が協調的に機能する場合や、配列の文脈依存的な効果を捉えることは困難である。

本研究の関連研究としてBranciamoreら（2018）\cite{branciamore2018intrinsic}の研究が挙げられる。この研究では\ac{tRNA}のアミノ酸決定因子をベイズネットワークによって推定している。問題点としてはtRNAdb\cite{abe2014trnadb}由来の整列済み配列（9758配列）を用い、アラインメントの縦位置を変数として扱うため、アラインメント依存の枠組みになっていること、「重要位置の推定・解釈」にあり、分類モデルとしての予測性能（外部テストでの汎化性能）を最適化・比較する設計ではない。
そのため、抽出された“重要位置”が、実際にアミノ酸分類などの下流タスク性能にどれだけ寄与するか、また未知データで再現するかは、別途検証が必要になることである。

\subsection{研究目的}

具体的に、以下の3つの研究課題を設定する。

\textbf{(1) 分類モデルの構築}

tRNA配列に対して \texttt{Domain\_AA} 形式のラベル（例：\texttt{bac\_Ala}）を付与し、RiNALMo\cite{penic2025rinalmo}に分類ヘッドを付加して多クラス分類としてファインチューニングを行い、分類モデルを構築した。

\textbf{(2) embedding解析}

層の学習推移を解析するため、embeddingに対するUMAP、中心コサイン類似度を用いた解析を行った。

\textbf{(3) \ac{attnlrp}\cite{achtibat2024attnlrp}による解析}

モデルの予測根拠を解釈可能にするため、逆伝播により分類スコアをrelevanceをトークンへ伝搬するAttnLRP実装を新規構築し、予測への寄与を算出することでモデルの解釈を試みた。

\clearpage

\section{材料}

研究に用いた具体的なソースコードに関しては\cref{subsec:source-code}、データセットに関しては\cref{subsec:datasets}に示す。

\subsection{ファインチューニングに用いるデータセット}

\subsubsection{tRNA配列データの取得}


本研究では、tRNA配列データとして主に3つの公開データベースから取得したデータを使用した。第一のデータソースはGtRNAdb\cite{chan2016gtrnadb}であり、tRNAscan-se\cite{chan2021trnascan}によって予測・同定されたtRNA遺伝子の包括的なコレクションを提供している。
第二のデータソースはtRNAdb-CE\cite{abe2014trnadb}であり、tRNAscan-SE\cite{chan2021trnascan}やARAgorn\cite{laslett2004aragorn}等のツールで予測したtRNAの中で信頼度が高いものをまとめている。
第三のデータソースはNCBI Datasets\cite{o2024exploring}である。
GtRNAdb\cite{chan2016gtrnadb}とtRNAdb-ce\cite{abe2014trnadb}ではArchaeaのデータが不足しているため、リファレンスゲノムから二つのデータベースに載っていないものを補完的に取得した。

データベースごとのドメインごとの件数については表\ref{tab:counts}に示す。

データベースからダウンロードした配列データはFASTA形式で提供されているが、RiNALMo\cite{penic2025rinalmo}への入力形式として、CSV形式でのデータ提供が要求されるため、FASTA形式からの変換処理を実施した。
CSVのカラムとしてsequence(tRNA配列)、amino\_acid(対応するアミノ酸)、domain(ドメイン)、anticodon(アンチコドン)となっている。抽出されたデータは、各カラムに対応する情報を持つCSVファイルとして保存され、後続の解析で使用された。

\begin{table}[H]
  \centering
  \begin{tabular}{lrrr}
    \hline
    データセット    & \ac{arc}         & \ac{bac}           & \ac{euk}         \\
    \hline
    GtRNAdb   & \numprint{5728}  & \numprint{48183}   & \numprint{6483}  \\
    tRNAdb-CE & \numprint{47092} & \numprint{9467233} & \numprint{78788} \\
    NCBI      & \numprint{32683} & 0                  & 0                \\
    % 合計 & 53820 & 9515416 & 176302 \\
    \hline
  \end{tabular}
  \caption{データセット別の件数}
  \label{tab:counts}
\end{table}



\subsubsection{データの前処理}

CSV形式への変換後、解析に適したデータセットを構築するために、複数段階の前処理を実施した。
以下の4段階に分かれる。

\noindent\textbf{(1) 配列の重複除去}\\
同一の配列に対して複数のラベルが付与されるとモデルの解釈性が低下すると考え、配列とラベルの対応関係をできるだけ一意に保つことを目的とした。
そのため、同一ラベル内で重複している同一配列は一つに統合し、異なるラベル間に同一配列が存在する場合は、該当する配列を双方のラベルから除去した。

\noindent\textbf{(2) 塩基による除外}\\
DNA配列の入力を統一し、学習および評価の安定性を高めることを目的として、前処理として4塩基（A, T, G, C）以外の文字（例：Nなど）を含む配列は解析対象から除外した。
これは、曖昧塩基を含む配列が入力のばらつきを増やし、学習や評価を不安定にする可能性があるためである。

\noindent\textbf{(3) tRNAscan-seによる検証}\\
元データの確認を行った際、明らかに配列長が長いものなどが存在したため、自らが収集したデータの正当性を検証するために、取得したデータに対してtRNAscan-se\cite{chan2021trnascan}を行った。
この際tRNAscan-se\cite{chan2021trnascan}のスコアが30より小さいものに関しては偽遺伝子の可能性があるため、論文を参考し除外している。
また元データと比較した際、対応するアミノ酸が異なっていたり、tRNAscan-se\cite{chan2021trnascan}が対応するアミノ酸を予測できない場合があったためそれらはデータベースに登録する際のエラーだと考え除外した。

以上基本的な前処理を行ったあとのデータセットについて表\ref{tab:aa_domain_counts}に示す。

\begin{table}[H]
  \centering
  {\small
    \setlength{\tabcolsep}{5pt}
    \begin{tabular}{l*{13}{r}}
      \toprule
      domain
       & Ala              & Arg              & Asn              & Asp              & Cys & Gln & Glu & Gly & His & Ile & Ile2 & Leu & Lys \\
      \midrule
      arc
       & \numprint{720}   & \numprint{1804}  & \numprint{297}   & \numprint{260}
       & \numprint{468}   & \numprint{541}   & \numprint{618}   & \numprint{914}
       & \numprint{312}   & \numprint{304}   & \numprint{4}     & \numprint{1698}
       & \numprint{594}                                                                                                                   \\
      bac
       & \numprint{11736} & \numprint{32933} & \numprint{7193}  & \numprint{5718}
       & \numprint{7371}  & \numprint{12708} & \numprint{10259} & \numprint{14674}
       & \numprint{8148}  & \numprint{5541}  & \numprint{51}    & \numprint{44042}
       & \numprint{12580}                                                                                                                 \\
      euk
       & \numprint{4454}  & \numprint{4619}  & \numprint{2074}  & \numprint{1924}
       & \numprint{1838}  & \numprint{2460}  & \numprint{2908}  & \numprint{3570}
       & \numprint{1329}  & \numprint{2369}  & \numprint{0}     & \numprint{3969}
       & \numprint{3251}                                                                                                                  \\
      \bottomrule
    \end{tabular}
  }

  \vspace{0.6em}

  {\small
    \setlength{\tabcolsep}{5.5pt}
    \begin{tabular}{l*{12}{r}}
      \toprule
      domain
       & Met              & Phe              & Pro              & Sec             & Ser & Sup & Thr & Trp & Tyr & Val & iMet & fMet \\
      \midrule
      arc
       & \numprint{570}   & \numprint{247}   & \numprint{749}   & \numprint{16}
       & \numprint{1668}  & \numprint{1}     & \numprint{1019}  & \numprint{172}
       & \numprint{259}   & \numprint{1029}  & \numprint{2}     & \numprint{0}                                                      \\
      bac
       & \numprint{16825} & \numprint{5295}  & \numprint{13566} & \numprint{2571}
       & \numprint{32119} & \numprint{82}    & \numprint{18608} & \numprint{6165}
       & \numprint{7599}  & \numprint{15278} & \numprint{0}     & \numprint{85}                                                     \\
      euk
       & \numprint{2008}  & \numprint{1621}  & \numprint{2705}  & \numprint{29}
       & \numprint{5705}  & \numprint{561}   & \numprint{3710}  & \numprint{1209}
       & \numprint{2905}  & \numprint{3713}  & \numprint{631}   & \numprint{0}                                                      \\
      \bottomrule
    \end{tabular}
  }
  \caption{基本的な前処理後のデータセット}
  \label{tab:aa_domain_counts}
\end{table}





\noindent\textbf{(4) 相同性除去}


訓練データとテストデータに高い相同性をもつ配列が同時に含まれると、モデルが本質的な特徴ではなく特定の配列パターンを記憶して過学習するおそれがある。
そこで前処理として、CD-HIT\cite{fu2012cd}を用いて配列間の相同性を低減した。

CD-HIT\cite{fu2012cd}では類似度の閾値を設定し、閾値以上で類似する配列を同一クラスタにまとめ、各クラスタから代表配列を1つだけ残すことで重複を除去する。
なお本データはラベル不均衡があるため、ドメインごとに-aS(短い方の配列長に対するアラインメント被覆率)を変更して（\ac{arc}: 0.95、\ac{euk}: 0.9、\ac{bac}: 0.8）クラスタリングした。
また、CD-HIT\cite{fu2012cd}はラベルごとに個別に実施している。

相同性除去後のデータセットを表\ref{tab:aa_domain_counts_cd-hit}に示す。

\begin{table}[H]
  \centering


  {\small
    \setlength{\tabcolsep}{5pt}
    \begin{tabular}{l*{13}{r}}
      \toprule
      domain
       & Ala             & Arg             & Asn             & Asp             & Cys & Gln & Glu & Gly & His & Ile & Ile2 & Leu & Lys \\
      \midrule
      arc
       & \numprint{325}  & \numprint{854}  & \numprint{162}  & \numprint{120}
       & \numprint{253}  & \numprint{285}  & \numprint{284}  & \numprint{447}
       & \numprint{157}  & \numprint{151}  & \numprint{3}    & \numprint{726}
       & \numprint{277}                                                                                                               \\
      bac
       & \numprint{1085} & \numprint{4111} & \numprint{823}  & \numprint{586}
       & \numprint{877}  & \numprint{1727} & \numprint{1219} & \numprint{1347}
       & \numprint{878}  & \numprint{937}  & \numprint{38}   & \numprint{4742}
       & \numprint{1310}                                                                                                              \\
      euk
       & \numprint{920}  & \numprint{1415} & \numprint{364}  & \numprint{378}
       & \numprint{355}  & \numprint{564}  & \numprint{653}  & \numprint{817}
       & \numprint{352}  & \numprint{510}  & \numprint{0}    & \numprint{967}
       & \numprint{667}                                                                                                               \\
      \bottomrule
    \end{tabular}
  }

  \vspace{0.6em}

  {\small
    \setlength{\tabcolsep}{5.5pt}
    \begin{tabular}{l*{12}{r}}
      \toprule
      domain
       & Met             & Phe             & Pro             & Sec            & Ser & Sup & Thr & Trp & Tyr & Val & iMet & fMet \\
      \midrule
      arc
       & \numprint{299}  & \numprint{133}  & \numprint{363}  & \numprint{9}
       & \numprint{716}  & \numprint{1}    & \numprint{494}  & \numprint{95}
       & \numprint{113}  & \numprint{464}  & \numprint{2}    & \numprint{0}                                                     \\
      bac
       & \numprint{1525} & \numprint{624}  & \numprint{1214} & \numprint{403}
       & \numprint{3151} & \numprint{61}   & \numprint{1797} & \numprint{907}
       & \numprint{742}  & \numprint{1495} & \numprint{0}    & \numprint{25}                                                    \\
      euk
       & \numprint{424}  & \numprint{324}  & \numprint{586}  & \numprint{10}
       & \numprint{1110} & \numprint{179}  & \numprint{953}  & \numprint{296}
       & \numprint{439}  & \numprint{917}  & \numprint{106}  & \numprint{0}                                                     \\
      \bottomrule
    \end{tabular}
  }
  \caption{相同性除去後のデータセット}
  \label{tab:aa_domain_counts_cd-hit}
\end{table}


\subsubsection{データセットの分割}

前処理後のデータセットを、train・val・testの3つに分割した。
これは、学習（train）、学習中の性能監視とモデル選択（val）、学習後の最終評価（test）を互いに独立なデータで行い、汎化性能を適切に測るためである。
なお、前処理後に分割が困難だった arc\_Sup,arc\_iMet,arc\_Ile2 は除外した。

分割に先立ち、ラベルを数値化した。
具体的には「ドメイン名\_アミノ酸」（例：bac\_Ala）を1つのラベルとして扱い、各ラベルに一意の整数（0,1,2, …）を割り当てて class\_id として付与し、元のラベルは name として保持した。
表\ref{tab:amino-acid-classid}にラベルと対応する class\_id を示し、最終的に68ラベルの分類モデルとして学習を行っている。

ドメインを含めたラベルにしたのは

(1) アミノ酸のみのラベルだとアンチコドン領域への過度な依存を招くと考えられること

(2) ドメインごとにtRNAの特徴が異なるため、ドメイン情報込みで学習させることで各ドメイン特有の配列特徴を捉えやすくするためである。

分割比率は train:val:test = 64:16:20 とした。
重要な点として、層化分割を実施した。具体的には、各データセットにおいてラベルの分布だけでなく、各ラベル内におけるアンチコドンの分布もできるだけ一致するように分割した。
本データセットでは、同一ラベル内でもアンチコドンの出現頻度に偏りが生じることがあり、これを考慮せずに分割すると、特定のアンチコドンが特定のサブセットに偏って評価が歪む可能性がある。
層化分割により、この偏りを抑え、より公平で安定した性能評価が可能となる。


\begin{table}[H]
  \centering
  \small
  \setlength{\tabcolsep}{8pt}
  \begin{tabular}{l r l r}
    \toprule
    amino\_acid & class\_id & amino\_acid & class\_id \\
    \midrule
    arc\_Ala    & 0         & bac\_Met    & 34        \\
    arc\_Arg    & 1         & bac\_Phe    & 35        \\
    arc\_Asn    & 2         & bac\_Pro    & 36        \\
    arc\_Asp    & 3         & bac\_Sec    & 37        \\
    arc\_Cys    & 4         & bac\_Ser    & 38        \\
    arc\_Gln    & 5         & bac\_Sup    & 39        \\
    arc\_Glu    & 6         & bac\_Thr    & 40        \\
    arc\_Gly    & 7         & bac\_Trp    & 41        \\
    arc\_His    & 8         & bac\_Tyr    & 42        \\
    arc\_Ile    & 9         & bac\_Val    & 43        \\
    arc\_Leu    & 10        & bac\_fMet   & 44        \\
    arc\_Lys    & 11        & euk\_Ala    & 45        \\
    arc\_Met    & 12        & euk\_Arg    & 46        \\
    arc\_Phe    & 13        & euk\_Asn    & 47        \\
    arc\_Pro    & 14        & euk\_Asp    & 48        \\
    arc\_Sec    & 15        & euk\_Cys    & 49        \\
    arc\_Ser    & 16        & euk\_Gln    & 50        \\
    arc\_Thr    & 17        & euk\_Glu    & 51        \\
    arc\_Trp    & 18        & euk\_Gly    & 52        \\
    arc\_Tyr    & 19        & euk\_His    & 53        \\
    arc\_Val    & 20        & euk\_Ile    & 54        \\
    bac\_Ala    & 21        & euk\_Leu    & 55        \\
    bac\_Arg    & 22        & euk\_Lys    & 56        \\
    bac\_Asn    & 23        & euk\_Met    & 57        \\
    bac\_Asp    & 24        & euk\_Phe    & 58        \\
    bac\_Cys    & 25        & euk\_Pro    & 59        \\
    bac\_Gln    & 26        & euk\_Sec    & 60        \\
    bac\_Glu    & 27        & euk\_Ser    & 61        \\
    bac\_Gly    & 28        & euk\_Sup    & 62        \\
    bac\_His    & 29        & euk\_Thr    & 63        \\
    bac\_Ile    & 30        & euk\_Trp    & 64        \\
    bac\_Ile2   & 31        & euk\_Tyr    & 65        \\
    bac\_Leu    & 32        & euk\_Val    & 66        \\
    bac\_Lys    & 33        & euk\_iMet   & 67        \\
    \bottomrule
  \end{tabular}
  \caption{amino\_acid と class\_id の対応表}
  \label{tab:amino-acid-classid}
\end{table}





\subsubsection{データセットの形状}
図\ref{fig:dataset_structure}に実際のデータセットのディレクトリ構造とCSVスキーマ（列定義）を示す。

\begin{figure}[H]
  \centering
  \fbox{
    \begin{minipage}{0.78\linewidth}
      \ttfamily
      dataset/\\
      \quad train.csv\\
      \quad val.csv\\
      \quad test.csv\\[2mm]
      CSV columns: name, sequence, class\_id
    \end{minipage}
  }
  \caption{データセットのディレクトリ構造とCSVスキーマ（列定義）．}
  \label{fig:dataset_structure}
\end{figure}


\subsection{内部表現解析に用いるデータセット}

内部表現解析として用いるデータセットとして、分類モデルにおいて正答と判断されたものを用いた。

また後続の解析においてデータ数が多いと解釈性が低下すること、また計算時間の効率化のため、各ラベルに対して36サンプルまでデータ数を絞った。
これはアミノ酸の同義コドン数の最小公倍数からこの数にしている。正答数が36に満たないものは除去してある。

この際ファインチューニング時と同じようにラベル内でアンチコドンが偏らないようにしており、その中でも予測確率\eqref{eq:Predicted_probability}が高いものを選択している。
ここで用いるラベルに関しては64個であり、表\ref{tab:label-set}に示す。

\begin{table}[H]
  \centering
  \small
  \setlength{\tabcolsep}{6pt}
  \begin{tabular}{llll}
    \toprule
    AA / special & \texttt{arc}      & \texttt{bac}      & \texttt{euk}       \\
    \midrule
    Ala          & \texttt{arc\_Ala} & \texttt{bac\_Ala} & \texttt{euk\_Ala}  \\
    Arg          & \texttt{arc\_Arg} & \texttt{bac\_Arg} & \texttt{euk\_Arg}  \\
    Asn          & \texttt{arc\_Asn} & \texttt{bac\_Asn} & \texttt{euk\_Asn}  \\
    Asp          & \texttt{arc\_Asp} & \texttt{bac\_Asp} & \texttt{euk\_Asp}  \\
    Cys          & \texttt{arc\_Cys} & \texttt{bac\_Cys} & \texttt{euk\_Cys}  \\
    Gln          & \texttt{arc\_Gln} & \texttt{bac\_Gln} & \texttt{euk\_Gln}  \\
    Glu          & \texttt{arc\_Glu} & \texttt{bac\_Glu} & \texttt{euk\_Glu}  \\
    Gly          & \texttt{arc\_Gly} & \texttt{bac\_Gly} & \texttt{euk\_Gly}  \\
    His          & \texttt{arc\_His} & \texttt{bac\_His} & \texttt{euk\_His}  \\
    Ile          & \texttt{arc\_Ile} & \texttt{bac\_Ile} & \texttt{euk\_Ile}  \\
    Leu          & \texttt{arc\_Leu} & \texttt{bac\_Leu} & \texttt{euk\_Leu}  \\
    Lys          & \texttt{arc\_Lys} & \texttt{bac\_Lys} & \texttt{euk\_Lys}  \\
    Met          & \texttt{arc\_Met} & \texttt{bac\_Met} & \texttt{euk\_Met}  \\
    Phe          & \texttt{arc\_Phe} & \texttt{bac\_Phe} & \texttt{euk\_Phe}  \\
    Pro          & \texttt{arc\_Pro} & \texttt{bac\_Pro} & \texttt{euk\_Pro}  \\
    Ser          & \texttt{arc\_Ser} & \texttt{bac\_Ser} & \texttt{euk\_Ser}  \\
    Thr          & \texttt{arc\_Thr} & \texttt{bac\_Thr} & \texttt{euk\_Thr}  \\
    Trp          & \texttt{arc\_Trp} & \texttt{bac\_Trp} & \texttt{euk\_Trp}  \\
    Tyr          & \texttt{arc\_Tyr} & \texttt{bac\_Tyr} & \texttt{euk\_Tyr}  \\
    Val          & \texttt{arc\_Val} & \texttt{bac\_Val} & \texttt{euk\_Val}  \\
    \midrule
    Sec          & --                & \texttt{bac\_Sec} & --                 \\
    Sup          & --                & \texttt{bac\_Sup} & \texttt{euk\_Sup}  \\
    iMet         & --                & --                & \texttt{euk\_iMet} \\
    \bottomrule
  \end{tabular}
  \caption{内部表現解析に用いるラベル一覧}
  \label{tab:label-set}
\end{table}

\clearpage

\section{方法}
本研究において用いたソースコードについては\cref{subsec:source-code}、パラメータは\cref{subsec:model-params}に示す。
\subsection{分類モデルの構築}

\subsubsection{RiNALMoの概要}
\label{sec:rinalmo_inst}

本研究では、RNA配列解析に特化した事前学習済みモデルとしてRiNALMo\cite{penic2025rinalmo}を採用した。
RiNALMo\cite{penic2025rinalmo}は、RNA配列を入力として自己教師あり学習により一般的な配列パターンや構造情報を内部表現として獲得する、BERT\cite{devlin2019bert}型のTransformerモデル\cite{vaswani2017attention}である。

RiNALMoの基本構造は、複数のTransformer blocksから構成され、各層は以下の主要構造を含む:

\begin{itemize}
  \item \textbf{\ac{MHA}}: 配列中の各トークンが他の全ての位置とどのように関連しているかを学習する。複数のAttentionヘッドにより、局所的・大域的な依存関係を並行して捉えていると考えられる。
  \item \textbf{\ac{FFN}}: 各トークンの表現を非線形変換により更新し、特徴の抽出と統合を行う。RiNALMo\cite{penic2025rinalmo}では近年の大規模モデルで用いられる活性化関数（SwiGLU\cite{shazeer2020glu}）を含む構成が採用されている。
  \item \textbf{add \& Norm}: 学習の安定性と勾配の伝播を改善するための正規化と残差接続である
\end{itemize}

入力トークン化は各ヌクレオチドを1トークンとして扱う方式としていて、モデル入力には系列の先頭に分類用トークン（[CLS]）を付与し、系列末尾には終端トークン（[EOS]）を付与する。
得られた最終層の出力から、配列全体の表現として[CLS]埋め込みを下流タスクの入力特徴として利用する。

位置情報の付与には\ac{Rope}\cite{su2024roformer}が用いられ、また大規模モデルの計算効率のためにFlashAttention-2\cite{dao2023flashattention}が採用されている。
これにより、長距離依存を含む文脈情報を効率的に表現へ取り込むことができる。

RiNALMo\cite{penic2025rinalmo}は、大規模なncRNA配列を用いた事前学習によって、RNA配列の一般的な「言語的パターン」を獲得している。
事前学習は\ac{MLM}に基づいて行われ、入力配列中の一部トークンをマスクし、周囲の文脈からマスクされたトークンを予測するタスクを通して学習される。
この自己教師あり学習により、モデルは配列の統計的規則性に加えて、構造・機能に関わる情報を内部表現として獲得していると考えられる。

\subsubsection{分類モデルの概要}
\label{sec:classification-model}
公開されている事前学習済み重みのgiga\cite{penic2025rinalmo}を初期値として利用した。gigaモデルは約650Mのパラメータ規模を持ち、埋め込み次元は1280、Transformer blockは33層から構成される。

また分類モデルの構築としては\ac{ncRNA}分類タスク\cite{penic2025rinalmo}が既に実装されているため、その分類ヘッドを再利用している。具体的には以下のような構成になっている。

分類ヘッドは2層MLPであり、
$\mathrm{Linear}(1280,256)\rightarrow \mathrm{GELU}\rightarrow \mathrm{Linear}(256,68)$
となっている。



\subsubsection{モデル全体の評価指標}
クラス数を $C$=68,評価データ数(test.csvの件数)を $N$=10332 とし、サンプル $i$ の正解ラベルを
$y_i\in\{0,\dots,C-1\}$，モデル出力ロジットを $\mathbf{z}_i\in\mathbb{R}^C$ とする。
Softmax により予測確率 $\mathbf{p}_i$ を
\begin{equation}
  p_{i,c}=\mathrm{softmax}(\mathbf{z}_i)_c
  =\frac{\exp(z_{i,c})}{\sum_{k=0}^{C-1}\exp(z_{i,k})}
  \quad (c=0,\dots,C-1)
  \label{eq:Predicted_probability}
\end{equation}
と定義する。以下に評価指標を示す。

\paragraph{(1) 損失（Cross-Entropy Loss）}
本研究では損失関数として交差エントロピー損失を用いており、
交差エントロピー損失は次式で与えられる：
\begin{equation}
  \mathcal{L}
  =-\frac{1}{N}\sum_{i=1}^{N}\log p_{i,y_i}.
  \label{eq:cross-entropy-loss}
\end{equation}

\paragraph{(2) 精度（Accuracy）}
予測ラベル $\hat{y}_i=\arg\max_{c}p_{i,c}$ を用いて，
\begin{equation}
  \mathrm{Accuracy}
  =\frac{1}{N}\sum_{i=1}^{N}\mathbf{1}\!\left[\hat{y}_i=y_i\right]
\end{equation}
と定義する。

\paragraph{(3) Weighted-F1}
クラス $c$ のsupport（真の件数）を $n_c=\sum_{i=1}^N \mathbf{1}[y_i=c]$ とし、
クラス別 $F1_c$(式\eqref{eq:F1}より) を support で重み付け平均して
\begin{equation}
  \mathrm{Weighted\mbox{-}F1}
  =\sum_{c=0}^{C-1}\frac{n_c}{N}\,F1_c
\end{equation}
とする。

\paragraph{(4) Macro-F1}
各クラスの $F1$ を等重みで平均し、
\begin{equation}
  \mathrm{Macro\mbox{-}F1}
  =\frac{1}{C}\sum_{c=0}^{C-1}F1_c
\end{equation}
とする。

\paragraph{(5) Micro-F1}
表\ref{tab:ovr_confusion}の定義に従って
全クラスで $\mathrm{TP},\mathrm{FP},\mathrm{FN}$ を合算すると以下のように定義される。
\begin{align}
  \mathrm{TP}_{\mu}               & =\sum_{c=0}^{C-1}\mathrm{TP}_c,\quad
  \mathrm{FP}_{\mu}=\sum_{c=0}^{C-1}\mathrm{FP}_c,\quad
  \mathrm{FN}_{\mu}=\sum_{c=0}^{C-1}\mathrm{FN}_c,                                                               \\
  \mathrm{Micro\mbox{-}Precision} & =\frac{\mathrm{TP}_{\mu}}{\mathrm{TP}_{\mu}+\mathrm{FP}_{\mu}},\quad
  \mathrm{Micro\mbox{-}Recall}=\frac{\mathrm{TP}_{\mu}}{\mathrm{TP}_{\mu}+\mathrm{FN}_{\mu}},                    \\
  \mathrm{Micro\mbox{-}F1}
                                  & =\frac{2\,\mathrm{Micro\mbox{-}Precision}\cdot \mathrm{Micro\mbox{-}Recall}}
  {\mathrm{Micro\mbox{-}Precision}+ \mathrm{Micro\mbox{-}Recall}}
\end{align}
と定義する。

\paragraph{(6) Top-3 Accuracy}
$\mathrm{Top3}(\mathbf{p}_i)$ を $\mathbf{p}_i$ の確率上位3クラスの集合とすると、
\begin{equation}
  \mathrm{Top\mbox{-}3\ Accuracy}
  =\frac{1}{N}\sum_{i=1}^{N}\mathbf{1}\!\left[y_i \in \mathrm{Top3}(\mathbf{p}_i)\right]
\end{equation}
と定義する。


\subsubsection{ラベル（クラス）ごとの評価指標}
各クラス $c$ について one-vs-rest とみなし，$\mathrm{TP}_c,\mathrm{FP}_c,\mathrm{FN}_c,\mathrm{TN}_c$
を定義する．混同行列（クラス $c$ に対する二値化）を表\ref{tab:ovr_confusion} に示す。

\begin{table}[H]
  \centering

  \label{tab:ovr_confusion}
  \begin{tabular}{c|cc}
    \hline
               & 予測：$c$          & 予測：$\neq c$     \\
    \hline
    真：$c$      & $\mathrm{TP}_c$ & $\mathrm{FN}_c$ \\
    真：$\neq c$ & $\mathrm{FP}_c$ & $\mathrm{TN}_c$ \\
    \hline
  \end{tabular}
  \caption{クラス $c$ に対する one-vs-rest の混同行列}
\end{table}

\paragraph{(1) 適合率（Precision）}
\begin{equation}
  \mathrm{Precision}_c
  =\frac{\mathrm{TP}_c}{\mathrm{TP}_c+\mathrm{FP}_c}.
  \label{eq:precison}
\end{equation}

\paragraph{(2) 再現率（Recall）}
\begin{equation}
  \mathrm{Recall}_c
  =\frac{\mathrm{TP}_c}{\mathrm{TP}_c+\mathrm{FN}_c}.
  \label{eq:recall}
\end{equation}

\paragraph{(3) F1}
\begin{equation}
  F1_c
  =\frac{2\,\mathrm{Precision}_c\cdot \mathrm{Recall}_c}
  {\mathrm{Precision}_c+\mathrm{Recall}_c}.
  \label{eq:F1}
\end{equation}

\subsubsection{誤分類Top10}
本当はクラスAなのに、クラスBと間違えたというペアをカウントし、最も頻度の高い誤りの上位10件を表示した。


\subsection{内部表現解析}
\ref{sec:classification-model}で示したように本研究では塩基配列をトークン列として扱い各塩基を1トークンに対応付ける。
入力系列長を $L$ とし、埋め込み次元を $d=1280$ とする。
系列の先頭には分類用の特殊トークン $\mathrm{[CLS]}$、末尾には終端トークン $\mathrm{[EOS]}$ を付与し、モデルへの入力は
\begin{equation}
  \mathbf{t}=\bigl(\mathrm{[CLS]},\, t_1,\, t_2,\, \dots,\, t_{L},\, \mathrm{[EOS]}\bigr)
\end{equation}
で表される（$t_i$ は $i$ 番目の塩基トークン）。

各トークンは各層で $d$ 次元ベクトルへ写像され，入力埋め込み行列
\begin{equation}
  \mathbf{E}=
  \begin{bmatrix}
    \mathbf{e}_{\mathrm{[CLS]}} \\
    \mathbf{e}_1                \\
    \vdots                      \\
    \mathbf{e}_{L}              \\
    \mathbf{e}_{\mathrm{[EOS]}}
  \end{bmatrix}
  \in \mathbb{R}^{(L+2)\times d}
  \label{eq:embed_matrix}
\end{equation}
を得る。ここで $\mathbf{e}_i\in\mathbb{R}^{d}$ は各塩基トークンに対応する埋め込みベクトルである。本モデルでは、Transformer本体はこの埋め込み $\mathbf{E}$ を入力として受け取り、自己注意機構を通じてトークン間の依存関係を表現する。

また、本研究の分類タスクでは最終層の $\mathrm{[CLS]}$ 表現 $\mathbf{h}_{\mathrm{[CLS]}}\in\mathbb{R}^{d}$ を系列全体の要約表現として用い、これを分類ヘッドへ入力することでロジットを算出している。

\subsubsection{UMAPによる可視化}
ここでは0,5,10,15,20,22,23,24,26,28,30,final層のUMAPを行う。

\textbf{(1)CLSトークンを用いたUMAP可視化}

層ごとにCLSトークンの埋め込みベクトルを抽出しDomain,AAごとにクラスタに分離されるか可視化した。
ここではDomain\_AA(モデルのラベル)に分けて可視化すると色分けが複雑になり可視性が低いため本論文には記載していない。

\textbf{(2)配列トークンを用いたUMAP可視化}

層ごとに配列トークンの埋め込みベクトルを抽出しDomain,AA,structure(tRNAの二次構造),Base(塩基)ごとにクラスタに分離されるか可視化した。
同様にDomain\_AAの結果は示さない。Domain,structureに関しては上に示した層の結果が乗っているが、Baseに関しては20層まで、AAに関しては22層以降結果を載せている。
\subsubsection{中心コサイン類似度解析}
中心コサイン類似度は以下の式で表すことができる。

% 次元 d = 1280, データ集合の平均ベクトルで中心化
\begin{equation}
  \operatorname{ccos}(x,y)
  =
  \frac{(x-\mu)^\top (y-\mu)}
  {\|x-\mu\|_2\,\|y-\mu\|_2}
  \qquad
  x,y\in\mathbb{R}^{1280}
  \label{eq:ccos}
\end{equation}

\begin{equation}
  \mu=\frac{1}{N}\sum_{n=1}^{N} z^{(n)}
  \qquad z^{(n)}\in\mathbb{R}^{1280}
  \label{eq:mu}
\end{equation}

ここで，$N$ はデータ数，$\mu$ はデータ集合の平均ベクトルである。
中心コサイン類似度は，データ集合の中心を基準にしたコサイン類似度を表し、データが中心からどの方向に分布しているかを捉えることができる。
値は-1から1の値をとり、1に近いほど類似していることを示す。
埋め込みベクトルの解析においては通常コサイン類似度が用いられるが、BERT型のモデルに関しては異方性の問題\cite{su2021whitening}があるとされることがあるため本研究では簡易的に中心コサイン類似度を用いている。

\textbf{(1)CLSトークンを用いた中心コサイン類似度解析}

層ごとにCLSトークンの埋め込みベクトルを抽出し、サンプルごとの中心コサイン類似度を計算した。
各塩基トークンに関しては計算量が大きいため、本研究では計算を行っていない。


\subsection{AttnLRPによる解析}

\subsubsection{AttnLRPの概要}


AttnLRP\cite{achtibat2024attnlrp}は、モデルの予測スコアから入力トークンへRelevanceを層ごとに逆向きに配分することで、最終予測に対する各入力要素の寄与を定量化する手法である。
本研究では、AttnLRP\cite{achtibat2024attnlrp}の公式githubリポジトリを参考にして、AttnLRPを本分類モデルに対して適用した。

まず，分類対象クラスのロジットを $s(\mathbf{x})$ とし，入力埋め込みを $\mathbf{E}\in\mathbb{R}^{（L+2)\times d}$（系列長 $L$，埋め込み次元 $d$）とする．トークン $i$ の寄与は
\begin{equation}
  R_i \;=\; \sum_{k=1}^{d} E_{i,k}\,\frac{\partial s}{\partial E_{i,k}}
  \label{eq:attnlrp_gxi}
\end{equation}
で与えられる。このとき、逆伝播される勾配がLRPの配分則に相当するよう、本研究で用いるAttnLRPでは各モジュールに以下のルールを適用している。

\paragraph{Attentionへの適用}
Scaled dot-product attention を
\begin{align}
  \mathbf{A} & = \mathrm{softmax}\!\left(\frac{\mathbf{Q}\mathbf{K}^{\top}}{\sqrt{d_h}}\right) \label{eq:A} \\
  \mathbf{C} & = \mathbf{A}\mathbf{V} \label{eq:C}
\end{align}
とおく（ヘッド次元 $d_h$）

AttnLRPでは行列積に対してrelevanceを一様に分配するために勾配を分割している。式\eqref{eq:C}より$\mathbf{A}$ と $\mathbf{V}$への勾配はそれぞれ1/2で分配される。更に式\eqref{eq:A}より$\mathbf{Q}$ と $\mathbf{K}$ への勾配も1/2で分配される。従って、$\mathbf{Q},\mathbf{K}$ には1/4ずつ、$\mathbf{V}$ には1/2の勾配が伝播されることになり、以下のように表すことができる($\mathrm{DG}(\cdot,\alpha)$ は逆伝播時に勾配を $1/\alpha$ にする演算)。
ただしsoftmaxの非線形性などが影響するため、厳密な保存則は成り立たないが、近似的に保存されると考える。
\begin{equation}
  \mathbf{Q}\leftarrow \mathrm{DG}(\mathbf{Q},4),\quad
  \mathbf{K}\leftarrow \mathrm{DG}(\mathbf{K},4),\quad
  \mathbf{V}\leftarrow \mathrm{DG}(\mathbf{V},2)
  \label{eq:attnlrp_divgrad}
\end{equation}


\paragraph{LayerNorm}
LayerNormは分散（標準偏差）計算を通る勾配がrelevance配分を歪めうるため、正規化の分母に対して勾配を遮断する：
\begin{equation}
  \mathbf{y}=\frac{\mathbf{x}-\mu(\mathbf{x})}{\mathrm{stopgrad}(\sigma(\mathbf{x}))}\,,
  \label{eq:attnlrp_layernorm}
\end{equation}
ここで $\mu,\sigma$ は特徴次元方向の平均・標準偏差である。これによりLayerNormを近似的に恒等写像として扱う。


\paragraph{Gated MLP（SwiGLU）}
SwiGLUを $ \mathbf{z}=f(\mathbf{u})\odot \mathbf{g}$（$f$ はSwish、$\odot$ は要素積）とおき、$\mathbf{u},\mathbf{g}$ をそれぞれ主枝側、ゲート側とすると、それぞれに勾配を分配すると考えられ、以下のように表すことができる。
\begin{equation}
  \mathbf{z}\leftarrow \mathrm{DG}(\mathbf{z},2)
  \label{eq:attnlrp_swiglu}
\end{equation}

以上のパッチにより，逆伝播で得られる $\partial s/\partial \mathbf{E}$ がAttnLRPの配分則を反映し、式\eqref{eq:attnlrp_gxi}によりトークン単位の寄与 $R_i$ を算出できる。

\subsubsection{AttnLRP実装に基づくRelevanceの算出手順}
実装は主に

(i) 各 Transformer ブロック内部の中間表現の取得（勾配保持）

(ii) 出力スコアに対する逆伝播

(iii) 活性値と勾配の内積に基づく Relevance の算出

(iv) Relevance 保存則を用いた残差寄与の復元の4段階からなる。以下に詳細を述べる。

\paragraph{(1) 中間出力の保存}

各 Transformer ブロック内のコンポーネントとしてはMHA,FFN,Residualが存在する。ここでforward 時に保存できるものは以下の3つの出力と勾配である。

\begin{itemize}
  \item \texttt{attn\_outputs}: Multi-Head Attention（MHA）の出力
  \item \texttt{ffn\_outputs}: Feed-Forward Network（FFN）の出力
  \item \texttt{block\_outputs}: 残差接続後のブロック全体出力（ブロック出力）
\end{itemize}


\paragraph{(2) 逆伝播（対象スコアに対する勾配計算）}
次に、モデルの予測スコア \(S\) に対して \texttt{backward()} を実行し、勾配を計算した。
分類設定における \(S\) は，通常の予測クラスのロジットとする。
また予測クラスのロジットとその他クラスのロジットの和との差を取ると他のクラスとの識別的な寄与を強調できるため、
（\(
S = \mathrm{logit}_{y^\ast}
- \sum_{y \neq y^\ast} \mathrm{logit}_{y}
\)
）
を用い、識別的な方向の寄与に焦点を当てており本研究ではこれを採用した。

\paragraph{(3) 活性値と勾配の内積によるRelevance算出}

Relevanceは、各中間表現の活性値と勾配の要素積の総和として算出した。
具体的に，あるコンポーネント出力 \(\mathbf{h}\) に対して
\begin{equation}
  R(\mathbf{h}) \;=\; \sum_{d} \left(h_d \cdot \frac{\partial S}{\partial h_d}\right)
  \label{eq:actgrad}
\end{equation}
と定義することができる（\(d\) は埋め込み次元）。

各 Transformer ブロック \(l\) において、以下の3種の Relevance を計算した：
\begin{align}
  R^{(l)}_{\mathrm{attn}}  & = \sum_d \left( \texttt{attn\_out}^{(l)}_d \cdot
  \left(\texttt{attn\_out}^{(l)}_d\right)'\right),                             \\
  R^{(l)}_{\mathrm{ffn}}   & = \sum_d \left( \texttt{ffn\_out}^{(l)}_d \cdot
  \left(\texttt{ffn\_out}^{(l)}_d\right)'\right),                              \\
  R^{(l)}_{\mathrm{total}} & = \sum_d \left( \texttt{block\_out}^{(l)}_d \cdot
  \left(\texttt{block\_out}^{(l)}_d\right)'\right),
\end{align}
ここで \((\cdot)'\) はスコア \(S\) に対する勾配（\(\partial S/\partial(\cdot)\)）を表す。

\paragraph{(4) Relevance保存則に基づく残差寄与の復元}
Transformer ブロックは残差接続を含むため、ブロック出力のRelevanceは
概念的には Attention,FFN,Residual の寄与に分解できる。
残差接続成分の Relevance を直接取得することはできないため、
ブロック全体の Relevance 保存則に基づき
\begin{equation}
  R^{(l)}_{\mathrm{res}} \;=\; R^{(l)}_{\mathrm{total}} - R^{(l)}_{\mathrm{attn}} - R^{(l)}_{\mathrm{ffn}}
  \label{eq:residual}
\end{equation}
として残差寄与を復元している。

\subsubsection{トークン特徴量ベクトルの構成}
得られた Relevance を用いてモデル内部の挙動を解析するため、各トークン \(t\) に対して全層の情報を集約した特徴量ベクトル \(\mathbf{x}_t\) を構成した。
具体的には，入力層における Residual 成分（入力埋め込みに対応）を \(r^{(0)}_{\mathrm{res}}\)、
第 \(l\) 層 (\(l=1, \dots, L\)) における Attention,FFN,Residual 各成分の Relevance をそれぞれ \(r^{(l)}_{\mathrm{attn}}, r^{(l)}_{\mathrm{ffn}}, r^{(l)}_{\mathrm{res}}\) とする。
本モデルの層数は \(L=33\) であるため、トークン \(t\) の特徴量ベクトル \(\mathbf{x}_t\) は以下の成分を連結したベクトルとして定義される。
\begin{equation}
  \mathbf{x}_t = \left[
    r^{(0)}_{\mathrm{res}}, \;
    \underbrace{r^{(1)}_{\mathrm{attn}}, r^{(1)}_{\mathrm{ffn}}, r^{(1)}_{\mathrm{res}}}_{\text{Layer 1}}, \;
    \dots, \;
    \underbrace{r^{(L)}_{\mathrm{attn}}, r^{(L)}_{\mathrm{ffn}}, r^{(L)}_{\mathrm{res}}}_{\text{Layer } L}
    \right]^\top
\end{equation}
このベクトルの次元数は
\begin{equation}
  1 + 3L = 1 + 3 \times 33 = 100
\end{equation}
となり、本研究では、この特徴量を用いて後段の解析および可視化を行う。

\subsubsection{層ごとのトークンタイプ別寄与率分析}
モデルの推論において、「どの層」で「どのトークンタイプ（CLS,Sequence Tokens,EOS）」が重要であるかを明らかにするため、トークンタイプごとの Relevance 集計を行った。
\paragraph{単一サンプルにおける計算}
ある入力配列について、トークン列を \(T\) とし、トークンタイプ集合を \(C = \{\texttt{CLS}, \texttt{Seq}, \texttt{EOS}\}\) とする。
タイプ \(c \in C\) に属するトークン集合を \(T_c \subset T\) と定義するとき、
第 \(l\) 層におけるタイプ \(c\) の Relevance の合計 \(R^{(l)}_c\) は次式で計算される。
\begin{equation}
  R^{(l)}_c = \sum_{t \in T_c} r^{(l)}_t
\end{equation}
ここで \(r^{(l)}_t\) は第 \(l\) 層におけるトークン \(t\) の Relevance（Attention、FFN、Residual 成分の和）である。
\paragraph{複数サンプルにわたる集計}
上記の計算を全サンプルにわたって集計する。
データセットに含まれる \(N\) 個のサンプルを \(s_1, s_2, \dots, s_N\) とし、
サンプル \(s\) における第 \(l\) 層・タイプ \(c\) の Relevance を \(R^{(l)}_c(s)\) と表記する。
\subparagraph{(1) 絶対値に基づく可視化（平均寄与量）}
各層・各トークンタイプが持つ Relevance の大きさを評価するため、全サンプルにわたる算術平均を算出した。
\begin{equation}
  \bar{R}^{(l)}_c = \frac{1}{N} \sum_{s=1}^{N} R^{(l)}_c(s)
\end{equation}
これにより、層が進むにつれて Relevance がどのように推移するかを絶対量として観察できる。
\subparagraph{(2) 比率に基づく可視化（平均寄与率）}
各層内での相対的な重要度を評価するため、まず各サンプルにおいて層内の全トークンタイプの Relevance に対する各タイプの割合を計算し、その後で全サンプルにわたる平均をとった。
サンプル \(s\) における第 \(l\) 層でのタイプ \(c\) の寄与率 \(P^{(l)}_c(s)\) は、
\begin{equation}
  P^{(l)}_c(s) = \frac{R^{(l)}_c(s)}{\sum_{c' \in C} R^{(l)}_{c'}(s)}
\end{equation}
であり、最終的な可視化には、この比率の全サンプル平均を用いた。
\begin{equation}
  \bar{P}^{(l)}_c = \frac{1}{N} \sum_{s=1}^{N} P^{(l)}_c(s)
\end{equation}
サンプルごとに比率を計算してから平均をとることで、Relevance の絶対値の大小（サンプルごとの差異）に影響されにくい、層ごとの構造的な寄与傾向を抽出した。

\subsubsection{層ごとのトークンタイプ内コンポーネント分析}
トークンタイプ別の分析をさらに詳細化し、各トークンタイプの Relevance が Attention、FFN、Residual のどのコンポーネントに由来するかを分析した。
\paragraph{単一サンプルにおける計算}
コンポーネント集合を \(K = \{\mathrm{attn}, \mathrm{ffn}, \mathrm{res}\}\) とする。
サンプル \(s\) の第 \(l\) 層において、トークンタイプ \(c\) のコンポーネント \(k\) からの Relevance の合計を \(R^{(l)}_{c,k}(s)\) と表記する。
トークンタイプ \(c\) の総 Relevance は、全コンポーネントの和として
\begin{equation}
  R^{(l)}_c(s) = \sum_{k \in K} R^{(l)}_{c,k}(s)
\end{equation}
で与えられる。
\paragraph{複数サンプルにわたる集計}
トークンタイプ別の分析と同様に、\(N\) 個のサンプルにわたって以下の2種類の集計を行った。
\subparagraph{(1) 絶対値に基づく可視化（平均寄与量）}
各コンポーネントの Relevance の大きさを全サンプルで平均した。
\begin{equation}
  \bar{R}^{(l)}_{c,k} = \frac{1}{N} \sum_{s=1}^{N} R^{(l)}_{c,k}(s)
\end{equation}
\subparagraph{(2) 比率に基づく可視化（符号付き寄与率）}
コンポーネント分析では、Relevance が正（促進的寄与）と負（抑制的寄与）の両方を取りうる。
単純な総和で正規化すると正負が相殺してしまうため、各コンポーネントの\textbf{絶対値の和}で正規化を行った。
これにより、符号（寄与の方向）を保持したまま、各コンポーネントの相対的な活動量を評価できる。
サンプル \(s\) における第 \(l\) 層・タイプ \(c\) のコンポーネント \(k\) の寄与率 \(Q^{(l)}_{c,k}(s)\) は次式で定義される。
\begin{equation}
  Q^{(l)}_{c,k}(s) = \frac{R^{(l)}_{c,k}(s)}{\sum_{k' \in K} \left| R^{(l)}_{c,k'}(s) \right|}
\end{equation}
分母は常に正であるため、\(Q^{(l)}_{c,k}(s)\) は元の Relevance と同じ符号を持つ。
正の大きな寄与は正の比率に、負の寄与は負の比率として表現される。
最終的な可視化には、この寄与率の全サンプル平均を用いた。
\begin{equation}
  \bar{Q}^{(l)}_{c,k} = \frac{1}{N} \sum_{s=1}^{N} Q^{(l)}_{c,k}(s)
\end{equation}


\subsubsection{単一サンプルにおけるrelevance解析}

各サンプルのRelevanceの流れと層ごとの変動を可視化・分析している。

\paragraph{可視化手法 (Relevance Flow Graph)}
Relevanceの推移とフローグラフとして描画している。
\begin{itemize}
  \item \textbf{ノード選択:} 可視化の複雑性を低減するため、各層においてRelevanceの絶対値が高い上位トークンから累積50\%のトークンのみを選択して表示している（ただしCLSトークンは常に含まれるようにしている）。
  \item \textbf{エッジ（情報の流れ）:} 層間のエッジはAttention機構による情報の移動を表している。層$L-1$のトークン$j$から層$L$のトークン$i$へのフロー量$F_{j \to i}$は、以下の式で定義する：
        \[
          F_{j \to i} = R_{L, i} \times A_{L, (i, j)}
        \]
        ここで、$R_{L, i}$は層$L$におけるトークン$i$のRelevance、$A_{L, (i, j)}$はAttention重みである。これは厳密なRelevanceの伝播ではないが、Attentionを介した情報の流れを近似的に表現している。
\end{itemize}

\paragraph{層間比較検定 (Layer Comparison Test)}
23層以降でのRelevanceの急激な変化に着目し、各塩基がモデルの最終判断においてどの程度重要視されるようになったかを統計的に評価するため、層22から最終層、input層から21層までのrelevanceの大きさを比較する検定を行っている。

各層でのRelevanceの大きさ（スケール）は異なる可能性があるため、まず各層内でRelevanceの標準化を行う。層$l$におけるトークン$i$のRelevanceを$R_{l, i}$とすると、標準化されたRelevance $Z_{l, i}$は次式で計算される。
\[
  Z_{l, i} = \frac{R_{l, i} - \mu_l}{\sigma_l}
\]
ここで、$\mu_l, \sigma_l$はそれぞれ層$l$における全トークンのRelevanceの平均と標準偏差である。

次に、層を「前期層（Layer 0-21）」と「後期層（Layer 22-Final）」の2つのグループに分け、各塩基トークン位置$i$について、後期層での標準化Relevanceが前期層よりも有意に高いかどうかを検証する。
これにはWelchのt検定（片側検定）を用いる。
\begin{itemize}
  \item 帰無仮説 $H_0$: $\mu_{late, i} \le \mu_{early, i}$
  \item 対立仮説 $H_1$: $\mu_{late, i} > \mu_{early, i}$
\end{itemize}
この検定により、後期層において、相対的な重要度が有意に上昇した塩基はモデルにとって重要な決定因子なのではないかという考察に基づき解析を行った。
この検定で示す結果には後期層においてz-scoreが正かつp値が0.05未満の塩基のみ示す。

すべてのサンプルを解析することは不可能なので、本研究ではarc\_Ala,bac\_Ala,euk\_Alaのラベルを持つサンプルを用いてAlaに関して解析を行っている。
ここで用いたサンプルは以下の通りである。またここで用いるサンプルに関してはtRNAscan-seを用いて、アノテーションを行っており、表\ref{fig:tRNAstructure}の塩基番号に合わせてある。

\begin{table}[H]
  \centering
  \begin{tabular}{l P{0.75\linewidth}}
    \hline
    label    & sequence                                                                                        \\
    \hline
    arc\_Ala & \seqsplit{GGGCTCGTAGATCAGTGGTAGATCGACACCTTCGCAAGGTGTAGGCCCCGGGTTCAAATCCCGGCGAGTCCA}             \\
    bac\_Ala & \seqsplit{GGGCTCGTAGATCAGTGGTAGATCGACACCTTCGCAAGGTGTAGGCCCCGGGTTCAAATCCCGGCGAGTCCA}             \\
    euk\_Ala & \seqsplit{GGGCGCTTGGTCTAGGGGTATGATCTTCGTTTCGCACCACACCGCGTCTACGAAAAGTCGTGGGTTCGATTCCTACAGCGTCCA} \\
    \hline
  \end{tabular}
  \caption{単一サンプルにおけるrelevance解析に用いるデータセット}
  \label{}
\end{table}


\clearpage
\section{結果}

\subsection{分類モデル}
\subsubsection{モデルの全体指標}

\begin{table}[H]
  \centering
  \begin{tabular}{lc}
    \toprule
    指標             & Test   \\
    \midrule
    Accuracy       & 0.9294 \\
    Macro-F1       & 0.8582 \\
    Micro-F1       & 0.9294 \\
    Weighted-F1    & 0.9286 \\
    Top-3 Accuracy & 0.9947 \\
    \bottomrule
  \end{tabular}
  \caption{RiNALMo分類モデルの評価結果（Test）}
  \label{tab:rinalmo_test_eval}
\end{table}

本研究では、クラス不均衡を含むtRNA配列データセットを用いて分類モデルを構築していることから、単一の評価指標に依存せず、複数の指標を用いてモデル性能を評価した。

AccuracyおよびMicro-F1はいずれも0.9264と高い値を示しており、モデルが全体として高精度な分類を実現していることが確認された。Weighted-F1も0.9246と同様に高い値を示しており、クラス分布を考慮した場合においても安定した性能を維持しているといえる。

一方で、Macro-F1は0.8582と他の指標と比較して低い値を示しており、これはサンプル数の少ないラベルにおいて識別が相対的に困難であることを反映していると考えられる。しかしながら、不均衡データであることを踏まえると、少数クラスを含めても0.85を超えるMacro-F1を達成している点は、本モデルが特定の多数クラスに過度に依存することなく、幅広いラベルの特徴を学習していることを示唆している。

加えて、Top-3 Accuracyは0.9947と極めて高い値を示しており、予測上位3クラスの中に正解アミノ酸が含まれる割合がほぼ100\%に近いことが確認された。

以上より、本モデルは不均衡なクラス分布下においても高い分類性能を有しており、後続の解析を行うための十分な性能を備えていると判断した。



\subsubsection{ラベルごとの評価}

\begin{table}[H]
  \centering


  \setlength{\tabcolsep}{3pt}
  \renewcommand{\arraystretch}{0.95}
  \footnotesize

  \begin{minipage}[t]{0.48\textwidth}
    \centering
    \begin{tabular}{lcccc}
      \toprule
      class\_name & Precision & Recall & F1     & Support \\
      \midrule
      arc\_Ala    & 0.9286    & 0.8000 & 0.8595 & 65      \\
      arc\_Arg    & 0.9310    & 0.9474 & 0.9391 & 171     \\
      arc\_Asn    & 0.7632    & 0.9062 & 0.8286 & 32      \\
      arc\_Asp    & 0.8696    & 0.8333 & 0.8511 & 24      \\
      arc\_Cys    & 0.9762    & 0.8039 & 0.8817 & 51      \\
      arc\_Gln    & 0.9483    & 0.9649 & 0.9565 & 57      \\
      arc\_Glu    & 0.8983    & 0.9464 & 0.9217 & 56      \\
      arc\_Gly    & 0.9043    & 0.9444 & 0.9239 & 90      \\
      arc\_His    & 0.9091    & 0.9677 & 0.9375 & 31      \\
      arc\_Ile    & 0.9000    & 0.9000 & 0.9000 & 30      \\
      arc\_Leu    & 0.8987    & 0.9793 & 0.9373 & 145     \\
      arc\_Lys    & 0.9811    & 0.9286 & 0.9541 & 56      \\
      arc\_Met    & 0.9298    & 0.8833 & 0.9060 & 60      \\
      arc\_Phe    & 0.9130    & 0.7778 & 0.8400 & 27      \\
      arc\_Pro    & 0.8667    & 0.8904 & 0.8784 & 73      \\
      arc\_Sec    & 0.0000    & 0.0000 & 0.0000 & 2       \\
      arc\_Ser    & 0.9338    & 0.8881 & 0.9104 & 143     \\
      arc\_Thr    & 0.9149    & 0.8687 & 0.8912 & 99      \\
      arc\_Trp    & 1.0000    & 0.9474 & 0.9730 & 19      \\
      arc\_Tyr    & 0.7333    & 0.9565 & 0.8302 & 23      \\
      arc\_Val    & 0.9444    & 0.9140 & 0.9290 & 93      \\
      bac\_Ala    & 0.9269    & 0.9312 & 0.9291 & 218     \\
      bac\_Arg    & 0.9500    & 0.9708 & 0.9603 & 822     \\
      bac\_Asn    & 0.9557    & 0.9152 & 0.9350 & 165     \\
      bac\_Asp    & 0.8462    & 0.9402 & 0.8907 & 117     \\
      bac\_Cys    & 0.9016    & 0.9886 & 0.9431 & 176     \\
      bac\_Gln    & 0.9768    & 0.9740 & 0.9754 & 346     \\
      bac\_Glu    & 0.9280    & 0.9508 & 0.9393 & 244     \\
      bac\_Gly    & 0.9011    & 0.9444 & 0.9222 & 270     \\
      bac\_His    & 0.9540    & 0.9486 & 0.9513 & 175     \\
      bac\_Ile    & 0.9184    & 0.9626 & 0.9399 & 187     \\
      bac\_Ile2   & 0.0000    & 0.0000 & 0.0000 & 7       \\
      bac\_Leu    & 0.9833    & 0.9325 & 0.9572 & 948     \\
      bac\_Lys    & 0.9341    & 0.9733 & 0.9533 & 262     \\
      bac\_Met    & 0.8690    & 0.9574 & 0.9111 & 305     \\
      \bottomrule
    \end{tabular}
  \end{minipage}\hfill
  \begin{minipage}[t]{0.48\textwidth}
    \centering
    \begin{tabular}{lcccc}
      \toprule
      class\_name & Precision & Recall & F1     & Support \\
      \midrule
      bac\_Phe    & 0.9273    & 0.8160 & 0.8681 & 125     \\
      bac\_Pro    & 0.9094    & 0.9545 & 0.9315 & 242     \\
      bac\_Sec    & 0.9634    & 0.9875 & 0.9753 & 80      \\
      bac\_Ser    & 0.9414    & 0.9683 & 0.9547 & 631     \\
      bac\_Sup    & 1.0000    & 0.6667 & 0.8000 & 12      \\
      bac\_Thr    & 0.9118    & 0.9499 & 0.9304 & 359     \\
      bac\_Trp    & 0.9778    & 0.9724 & 0.9751 & 181     \\
      bac\_Tyr    & 0.9854    & 0.9122 & 0.9474 & 148     \\
      bac\_Val    & 0.9559    & 0.9431 & 0.9495 & 299     \\
      bac\_fMet   & 0.0000    & 0.0000 & 0.0000 & 5       \\
      euk\_Ala    & 0.9319    & 0.9674 & 0.9493 & 184     \\
      euk\_Arg    & 0.9286    & 0.8697 & 0.8982 & 284     \\
      euk\_Asn    & 0.9054    & 0.9306 & 0.9178 & 72      \\
      euk\_Asp    & 0.9365    & 0.7867 & 0.8551 & 75      \\
      euk\_Cys    & 0.9836    & 0.8451 & 0.9091 & 71      \\
      euk\_Gln    & 0.9545    & 0.9292 & 0.9417 & 113     \\
      euk\_Glu    & 0.9583    & 0.8779 & 0.9163 & 131     \\
      euk\_Gly    & 0.9400    & 0.8598 & 0.8981 & 164     \\
      euk\_His    & 0.9118    & 0.8732 & 0.8921 & 71      \\
      euk\_Ile    & 0.9574    & 0.8824 & 0.9184 & 102     \\
      euk\_Leu    & 0.7754    & 0.9482 & 0.8531 & 193     \\
      euk\_Lys    & 0.9520    & 0.8947 & 0.9225 & 133     \\
      euk\_Met    & 0.8507    & 0.6706 & 0.7500 & 85      \\
      euk\_Phe    & 0.7468    & 0.9219 & 0.8252 & 64      \\
      euk\_Pro    & 0.9709    & 0.8547 & 0.9091 & 117     \\
      euk\_Sec    & 0.0000    & 0.0000 & 0.0000 & 2       \\
      euk\_Ser    & 0.9565    & 0.8919 & 0.9231 & 222     \\
      euk\_Sup    & 0.7692    & 0.8333 & 0.8000 & 36      \\
      euk\_Thr    & 0.9185    & 0.8848 & 0.9013 & 191     \\
      euk\_Trp    & 0.9048    & 0.9661 & 0.9344 & 59      \\
      euk\_Tyr    & 0.9341    & 0.9659 & 0.9497 & 88      \\
      euk\_Val    & 0.8974    & 0.9563 & 0.9259 & 183     \\
      euk\_iMet   & 1.0000    & 0.9524 & 0.9756 & 21      \\
      \bottomrule
    \end{tabular}
  \end{minipage}
  \caption{RiNALMo分類モデルのラベル別評価指標（Test）}
  \label{tab:rinalmo_test_per_label}
\end{table}


表\ref{tab:rinalmo_test_per_label}より、ラベルごとの性能には明確な差が確認された。まず、多数のラベルでF1が0.9前後に達しており、特にDomainがBacteriaのラベルでは高水準で安定している。
例えば、bac\_Leu（F1=0.9572, Support=948）、bac\_Gln（F1=0.9754, Support=346）、bac\_Arg（F1=0.9603, Support=822）、bac\_Ser（F1=0.9547, Support=631）などはPrecision・Recallともに高く、十分なsupportが存在するラベルに関しては高い識別性能を示している。

一方で、少数クラスや特殊クラスでは性能低下が顕著である。
特に、bac\_Ile2（Support=7）、bac\_fMet（Support=5）、arc\_Sec（Support=2）、euk\_Sec（Support=2）はPrecision/Recall/F1がすべて0となっており、モデルが当該クラスを一度も正しく予測できていない。

\subsubsection{誤分類Top10}

\begin{table}[H]
  \centering

  \begin{tabular}{llr}
    \toprule
    True label & Predicted label & Count \\
    \midrule
    bac\_Leu   & euk\_Leu        & 49    \\
    euk\_Arg   & bac\_Arg        & 34    \\
    euk\_Met   & bac\_Met        & 27    \\
    euk\_Gly   & bac\_Gly        & 22    \\
    euk\_Thr   & bac\_Thr        & 19    \\
    euk\_Ser   & bac\_Ser        & 19    \\
    bac\_Phe   & euk\_Phe        & 18    \\
    arc\_Ser   & bac\_Ser        & 16    \\
    euk\_Pro   & bac\_Pro        & 16    \\
    euk\_Asp   & bac\_Asp        & 16    \\
    \bottomrule
  \end{tabular}
  \caption{誤分類ペア上位10件（Test）}
  \label{tab:misclf_top10}
\end{table}

表\ref{tab:misclf_top10}では同一アミノ酸を維持したままドメインのみが置換されるパターンが大半を占めていて、特にEukaryotaからBacteriaへの誤分類が顕著である。



\subsection{内部表現解析}

\subsubsection{中心コサイン類似度分析}
\paragraph{ヒートマップの構造とカラーマップ}
\begin{figure}[H]
  \centering
  \includegraphics[width=0.60\textwidth]{figures/CCS_heatmaps_cls_thesis/similarity_colormap.png}
  \caption{Centered Cosine Similarityのカラーマップ}
  \label{fig:similarity_colormap}
\end{figure}

表\ref{tab:label-set}に示すように、各Domain\_AAクラスは64でそれぞれのサンプル数は36である。
全 \(N = 36\times 64 = 2304\) サンプルについて \(N \times N \) の類似度行列を計算し、ヒートマップとして可視化した。
サンプルはDomain\_AAのアルファベット順にソートして並べた。
図~\ref{fig:heatmap_structure} にヒートマップの構造を示す。
\begin{figure}[htbp]
  \centering
  \begin{tikzpicture}[scale=0.6]
    % Viridis カラー定義（0→1 : 紫→青→緑→黄）
    \definecolor{viridis0}{RGB}{68, 1, 84}      % 低（紫）
    \definecolor{viridis3}{RGB}{59, 82, 139}   % 中低（青紫）
    \definecolor{viridis5}{RGB}{33, 145, 140}  % 中（ティール）
    \definecolor{viridis7}{RGB}{94, 201, 98}   % 中高（黄緑）
    \definecolor{viridis10}{RGB}{253, 231, 37} % 高（黄）

    % グリッドサイズ（6x6 ブロック、各ブロック2x2単位）
    \def\blocksize{2}
    \def\numblocks{6}
    \def\totalsize{12}

    % オフダイアゴナルの塗りつぶし（低〜中類似度）
    \foreach \i in {0,...,5} {
        \foreach \j in {0,...,5} {
            \pgfmathsetmacro{\xstart}{\i * \blocksize}
            \pgfmathsetmacro{\ystart}{(\numblocks - 1 - \j) * \blocksize}
            \ifnum\i=\j
              % 対角ブロック：高類似度（黄色）
              \fill[viridis10] (\xstart, \ystart) rectangle +(\blocksize, \blocksize);
            \else
              % オフダイアゴナル：低〜中類似度（ティール〜青紫）
              \fill[viridis3] (\xstart, \ystart) rectangle +(\blocksize, \blocksize);
            \fi
          }
      }

    % 境界線（白）
    \foreach \i in {1,...,5} {
        \draw[white, line width=0.8pt] (\i*\blocksize, 0) -- (\i*\blocksize, \totalsize);
        \draw[white, line width=0.8pt] (0, \i*\blocksize) -- (\totalsize, \i*\blocksize);
      }

    % 外枠
    \draw[thick] (0,0) rectangle (\totalsize, \totalsize);

    % アミノ酸ラベル（例：6種）
    \def\labels{{"arc\_Ala", "arc\_Arg", "arc\_Asn", "arc\_Asp", "arc\_Cys", "arc\_Gln"}}

    % 左側ラベル（Y軸）
    \foreach \i in {0,...,5} {
        \pgfmathsetmacro{\ypos}{(\numblocks - 0.5 - \i) * \blocksize}
        \pgfmathsetmacro{\labelidx}{int(\i)}
        \node[left, font=\small] at (-0.3, \ypos) {\pgfmathparse{\labels[\labelidx]}\pgfmathresult};
      }

    % 下側ラベル（X軸）
    \foreach \i in {0,...,5} {
        \pgfmathsetmacro{\xpos}{(\i + 0.5) * \blocksize}
        \pgfmathsetmacro{\labelidx}{int(\i)}
        \node[below, rotate=45, anchor=north east, font=\small] at (\xpos, -0.3) {\pgfmathparse{\labels[\labelidx]}\pgfmathresult};
      }

    % 注釈矢印
    \draw[<-, thick] (1, 11) -- (14, 11) node[right, align=left, font=\small] {対角ブロック\\同一Domain\_AA間};
    \draw[<-, thick] (5, 9) -- (14, 9) node[right, align=left, font=\small] {非対角ブロック\\異なるDomain\_AA間};

  \end{tikzpicture}
  \caption[CLSトークン類似度ヒートマップの構造]{
    CLS トークン類似度ヒートマップの構造。
    サンプルはDomain\_AAのアルファベット順に並べられている。
    各アミノ酸のサンプル数は均等であり、ブロックサイズも同一である。
    対角ブロック（黄色）は同一Domain\_AAを運ぶ tRNA 配列間の類似度を、
    非対角ブロック（青紫）は異なるDomain\_AA間の類似度を表す。
    カラーマップは Viridis を使用し、値域は \([0, 1]\) である。これは今回の解析では類似度が負の値をとることがなかったためである。
  }
  \label{fig:heatmap_structure}
\end{figure}
このヒートマップにより、以下を読み取ることができる：
\begin{itemize}
  \item \textbf{対角ブロック}：同一Domain\_AAを運ぶ tRNA 配列同士の類似度。高い類似度（明るい色）であれば、モデルが同一Domain\_AAの配列を類似した表現で捉えていることを示す。
  \item \textbf{非対角ブロック}：異なるDomain\_AA間の類似度。特定のDomain\_AAペア間で高い類似度が見られれば、モデルがそれらを類似したカテゴリとして認識している可能性がある。
\end{itemize}

\subsubsection{UMAPを用いた可視化}
\begin{figure}[H]
  \centering
  \begin{subfigure}{0.38\linewidth}
    \centering
    \includegraphics[width=\linewidth,height=0.25\textheight,keepaspectratio]{figures/umap_cls_thesis/AA_legend.png}
    \caption{CLSのUMAP可視化用AAカラーマップ}
  \end{subfigure}
  \hfill
  \begin{subfigure}{0.38\linewidth}
    \centering
    \includegraphics[width=\linewidth,height=0.25\textheight,keepaspectratio]{figures/umap_cls_thesis/Domain_legend.png}
    \caption{CLSのUMAP可視化用Domainカラーマップ}
  \end{subfigure}
  \caption{CLSのUMAP可視化用カラーマップ}
\end{figure}

\begin{figure}[H]
  \centering
  \begin{subfigure}{0.38\linewidth}
    \centering
    \includegraphics[width=\linewidth,height=0.25\textheight,keepaspectratio]{figures/umap_token_thesis/legends/legend_AA.png}
    \caption{各塩基トークンのUMAP可視化用AAカラーマップ}
  \end{subfigure}
  \hfill
  \begin{subfigure}{0.38\linewidth}
    \centering
    \includegraphics[width=\linewidth,height=0.25\textheight,keepaspectratio]{figures/umap_token_thesis/legends/legend_Domain.png}
    \caption{各塩基トークンのUMAP可視化用Domainカラーマップ}
  \end{subfigure}
  \caption{各塩基トークンのUMAP可視化用カラーマップ(AA,Domain)}
\end{figure}

\begin{figure}[H]
  \centering
  \begin{subfigure}{0.38\linewidth}
    \centering
    \includegraphics[width=\linewidth,height=0.25\textheight,keepaspectratio]{figures/umap_token_thesis_BASE/legends/legend_Base.png}
    \caption{各塩基トークンのUMAP可視化用Baseカラーマップ}
  \end{subfigure}
  \hfill
  \begin{subfigure}{0.38\linewidth}
    \centering
    \includegraphics[width=\linewidth,height=0.25\textheight,keepaspectratio]{figures/umap_token_thesis/legends/legend_Structure.png}
    \caption{各塩基トークンのUMAP可視化用Structureカラーマップ}
  \end{subfigure}
  \caption{各塩基トークンのUMAP可視化用カラーマップ(AA,Domain)}
\end{figure}

\begin{table}[H]
  \centering
  \begin{tabular}{llp{5cm}}
    \hline
    \textbf{ラベル}                                                   & \textbf{部位}      & \textbf{区分}                \\
    \hline
    \texttt{acceptor\_stem\_5side} / \texttt{\_3side}              & Acceptor arm     & ステム（5'側鎖 / 3'側鎖）           \\
    \texttt{acceptor\_D\_linker}                                   & Acceptor--D間     & acceptor arm,D armの接続部     \\
    \hline
    \texttt{D\_stem\_5side} / \texttt{\_3side}                     & D arm            & ステム（5'側鎖 / 3'側鎖）           \\
    \texttt{D\_loop}                                               & D arm            & ループ                        \\
    \texttt{D\_anticodon\_linker}                                  & D--Anticodon間    & D arm,Anticodon armの接続部    \\
    \hline
    \texttt{anticodon\_stem\_5side} / \texttt{\_3side}             & Anticodon arm    & ステム（5'側鎖 / 3'側鎖）           \\
    \texttt{anticodon\_loop\_5side} / \texttt{\_3side}             & Anticodon arm    & ループ（5'側 / 3'側）             \\
    \texttt{anticodon\_core}                                       & Anticodon arm    & アンチコドン                     \\
    \texttt{anticodon\_VR\_linker} / \texttt{anticodon\_T\_linker} & Anticodon--VR/T間 & Anticodon armとVR/T armの接続部 \\
    \hline
    \texttt{VR\_stem\_5side} / \texttt{\_3side}                    & VR               & ステム（5'側鎖 / 3'側鎖）           \\
    \texttt{VR\_loop}                                              & VR               & ループ（loop）                  \\
    \texttt{VR\_T\_linker}                                         & VR--T間           & VRとT armの接続部               \\
    \hline
    \texttt{T\_stem\_5side} / \texttt{\_3side}                     & T arm            & ステム（5'側鎖 / 3'側鎖）           \\
    \texttt{T\_loop}                                               & T arm            & ループ（T$\Psi$Cループ）           \\
    \texttt{T\_acceptor\_linker}                                   & T--Acceptor間     & T armとAcceptor armの接続部     \\
    \hline
    \texttt{leader\_5side}                                         & 5' leader        & pre-tRNAの5'側余分配列（成熟で除去）    \\
    \texttt{leader\_3side}                                         & 3' trailer       & pre-tRNAの3'側余分配列（成熟で除去）    \\
    \hline
  \end{tabular}
  \caption{structureとtRNA二次構造の対応}
  \label{tab:trna_structure_labels}
\end{table}


\subsubsection{CLSの内部表現解析}


\begin{figure}[H]
  \centering

  % 1行目
  \begin{subfigure}[t]{0.48\linewidth}
    \centering
    \includegraphics[width=\linewidth]{figures/CCS_heatmaps_cls_thesis/layer_00/layer_00_heatmap.png}
    \caption{layer 0：CLS：中心コサイン類似度ヒートマップ}
  \end{subfigure}\hfill
  \begin{subfigure}[t]{0.48\linewidth}
    \centering
    \includegraphics[width=\linewidth]{figures/CCS_heatmaps_cls_thesis/layer_05/layer_05_heatmap.png}
    \caption{layer 5：CLS：中心コサイン類似度ヒートマップ}
  \end{subfigure}

  \vspace{0.6em}

  % 2行目
  \begin{subfigure}[t]{0.48\linewidth}
    \centering
    \includegraphics[width=\linewidth]{figures/umap_cls_thesis/layer_00/umap_cls_by_aa.png}
    \caption{layer 0：CLS：UMAP可視化（AA）}
  \end{subfigure}\hfill
  \begin{subfigure}[t]{0.48\linewidth}
    \centering
    \includegraphics[width=\linewidth]{figures/umap_cls_thesis/layer_05/umap_cls_by_aa.png}
    \caption{layer 5：CLS：UMAP可視化（AA）}
  \end{subfigure}

  \vspace{0.6em}

  % 3行目
  \begin{subfigure}[t]{0.48\linewidth}
    \centering
    \includegraphics[width=\linewidth]{figures/umap_cls_thesis/layer_00/umap_cls_by_domain.png}
    \caption{layer 0：CLS：UMAP可視化（Domain）}
  \end{subfigure}\hfill
  \begin{subfigure}[t]{0.48\linewidth}
    \centering
    \includegraphics[width=\linewidth]{figures/umap_cls_thesis/layer_05/umap_cls_by_domain.png}
    \caption{layer 5：CLS：UMAP可視化（Domain）}
  \end{subfigure}

  \caption{layer0とlayer5のCLSトークンに関する中心コサイン類似度ヒートマップおよびUMAP可視化}
  \label{fig:cls0_5}
\end{figure}


\begin{figure}[H]
  \centering

  % 1行目
  \begin{subfigure}[t]{0.48\linewidth}
    \centering
    \includegraphics[width=\linewidth]{figures/CCS_heatmaps_cls_thesis/layer_10/layer_10_heatmap.png}
    \caption{layer 10：CLS：中心コサイン類似度ヒートマップ}
  \end{subfigure}\hfill
  \begin{subfigure}[t]{0.48\linewidth}
    \centering
    \includegraphics[width=\linewidth]{figures/CCS_heatmaps_cls_thesis/layer_15/layer_15_heatmap.png}
    \caption{layer 15：CLS：中心コサイン類似度ヒートマップ}
  \end{subfigure}

  \vspace{0.6em}

  % 2行目
  \begin{subfigure}[t]{0.48\linewidth}
    \centering
    \includegraphics[width=\linewidth]{figures/umap_cls_thesis/layer_10/umap_cls_by_aa.png}
    \caption{layer 10：CLS：UMAP可視化（AA）}
  \end{subfigure}\hfill
  \begin{subfigure}[t]{0.48\linewidth}
    \centering
    \includegraphics[width=\linewidth]{figures/umap_cls_thesis/layer_15/umap_cls_by_aa.png}
    \caption{layer 15：CLS：UMAP可視化（AA）}
  \end{subfigure}

  \vspace{0.6em}

  % 3行目
  \begin{subfigure}[t]{0.48\linewidth}
    \centering
    \includegraphics[width=\linewidth]{figures/umap_cls_thesis/layer_10/umap_cls_by_domain.png}
    \caption{layer 10：CLS：UMAP可視化（Domain）}
  \end{subfigure}\hfill
  \begin{subfigure}[t]{0.48\linewidth}
    \centering
    \includegraphics[width=\linewidth]{figures/umap_cls_thesis/layer_15/umap_cls_by_domain.png}
    \caption{layer 15：CLS：UMAP可視化（Domain）}
  \end{subfigure}

  \caption{layer10とlayer15のCLSトークンに関する中心コサイン類似度ヒートマップおよびUMAP可視化}
  \label{fig:cls10_15}
\end{figure}

\begin{figure}[H]
  \centering

  % 1行目
  \begin{subfigure}[t]{0.48\linewidth}
    \centering
    \includegraphics[width=\linewidth]{figures/layer_20_heatmap_abst.png}
    \caption{layer 20：CLS：中心コサイン類似度ヒートマップ}
  \end{subfigure}\hfill
  \begin{subfigure}[t]{0.48\linewidth}
    \centering
    \includegraphics[width=\linewidth]{figures/layer_22_heatmap_abst.png}
    \caption{layer 22：CLS：中心コサイン類似度ヒートマップ}
  \end{subfigure}

  \vspace{0.6em}

  % 2行目
  \begin{subfigure}[t]{0.48\linewidth}
    \centering
    \includegraphics[width=\linewidth]{figures/umap_cls_thesis/layer_20/umap_cls_by_aa.png}
    \caption{layer 20：CLS：UMAP可視化（AA）}
  \end{subfigure}\hfill
  \begin{subfigure}[t]{0.48\linewidth}
    \centering
    \includegraphics[width=\linewidth]{figures/umap_cls_thesis/layer_22/umap_cls_by_aa.png}
    \caption{layer 22：CLS：UMAP可視化（AA）}
  \end{subfigure}

  \vspace{0.6em}

  % 3行目
  \begin{subfigure}[t]{0.48\linewidth}
    \centering
    \includegraphics[width=\linewidth]{figures/umap_cls_thesis/layer_20/umap_cls_by_domain.png}
    \caption{layer 20：CLS：UMAP可視化（Domain）}
  \end{subfigure}\hfill
  \begin{subfigure}[t]{0.48\linewidth}
    \centering
    \includegraphics[width=\linewidth]{figures/umap_cls_thesis/layer_22/umap_cls_by_domain.png}
    \caption{layer 22：CLS：UMAP可視化（Domain）}
  \end{subfigure}

  \caption{layer20とlayer22のCLSトークンに関する中心コサイン類似度ヒートマップおよびUMAP可視化}
  \label{fig:cls20_22}
\end{figure}

\begin{figure}[H]
  \centering

  % 1行目
  \begin{subfigure}[t]{0.48\linewidth}
    \centering
    \includegraphics[width=\linewidth]{figures/CCS_heatmaps_cls_thesis_red/layer_23/layer_23_heatmap.png}
    \caption{layer 23：CLS：中心コサイン類似度ヒートマップ}
  \end{subfigure}\hfill
  \begin{subfigure}[t]{0.48\linewidth}
    \centering
    \includegraphics[width=\linewidth]{figures/CCS_heatmaps_cls_thesis_red/layer_24/layer_24_heatmap.png}
    \caption{layer 24：CLS：中心コサイン類似度ヒートマップ}
  \end{subfigure}

  \vspace{0.6em}

  % 2行目
  \begin{subfigure}[t]{0.48\linewidth}
    \centering
    \includegraphics[width=\linewidth]{figures/umap_cls_thesis/layer_23/umap_cls_by_aa.png}
    \caption{layer 23：CLS：UMAP可視化（AA）}
  \end{subfigure}\hfill
  \begin{subfigure}[t]{0.48\linewidth}
    \centering
    \includegraphics[width=\linewidth]{figures/umap_cls_thesis/layer_24/umap_cls_by_aa.png}
    \caption{layer 24：CLS：UMAP可視化（AA）}
  \end{subfigure}

  \vspace{0.6em}

  % 3行目
  \begin{subfigure}[t]{0.48\linewidth}
    \centering
    \includegraphics[width=\linewidth]{figures/umap_cls_thesis/layer_23/umap_cls_by_domain.png}
    \caption{layer 23：CLS：UMAP可視化（Domain）}
  \end{subfigure}\hfill
  \begin{subfigure}[t]{0.48\linewidth}
    \centering
    \includegraphics[width=\linewidth]{figures/umap_cls_thesis/layer_24/umap_cls_by_domain.png}
    \caption{layer 24：CLS：UMAP可視化（Domain）}
  \end{subfigure}

  \caption{layer23とlayer24のCLSトークンに関する中心コサイン類似度ヒートマップおよびUMAP可視化}
  \label{fig:cls23_24}
\end{figure}


\begin{figure}[H]
  \centering

  % 1行目
  \begin{subfigure}[t]{0.48\linewidth}
    \centering
    \includegraphics[width=\linewidth]{figures/CCS_heatmaps_cls_thesis_red/layer_26/layer_26_heatmap.png}
    \caption{layer 26：CLS：中心コサイン類似度ヒートマップ}
  \end{subfigure}\hfill
  \begin{subfigure}[t]{0.48\linewidth}
    \centering
    \includegraphics[width=\linewidth]{figures/CCS_heatmaps_cls_thesis_red/layer_28/layer_28_heatmap.png}
    \caption{layer 28：CLS：中心コサイン類似度ヒートマップ}
  \end{subfigure}

  \vspace{0.6em}

  % 2行目
  \begin{subfigure}[t]{0.48\linewidth}
    \centering
    \includegraphics[width=\linewidth]{figures/umap_cls_thesis/layer_26/umap_cls_by_aa.png}
    \caption{layer 26：CLS：UMAP可視化（AA）}
  \end{subfigure}\hfill
  \begin{subfigure}[t]{0.48\linewidth}
    \centering
    \includegraphics[width=\linewidth]{figures/umap_cls_thesis/layer_28/umap_cls_by_aa.png}
    \caption{layer 28：CLS：UMAP可視化（AA）}
  \end{subfigure}

  \vspace{0.6em}

  % 3行目
  \begin{subfigure}[t]{0.48\linewidth}
    \centering
    \includegraphics[width=\linewidth]{figures/umap_cls_thesis/layer_26/umap_cls_by_domain.png}
    \caption{layer 26：CLS：UMAP可視化（Domain）}
  \end{subfigure}\hfill
  \begin{subfigure}[t]{0.48\linewidth}
    \centering
    \includegraphics[width=\linewidth]{figures/umap_cls_thesis/layer_28/umap_cls_by_domain.png}
    \caption{layer 28：CLS：UMAP可視化（Domain）}
  \end{subfigure}

  \caption{layer26とlayer28のCLSトークンに関する中心コサイン類似度ヒートマップおよびUMAP可視化}
  \label{fig:cls26_28}
\end{figure}

\begin{figure}[H]
  \centering

  % 1行目
  \begin{subfigure}[t]{0.48\linewidth}
    \centering
    \includegraphics[width=\linewidth]{figures/CCS_heatmaps_cls_thesis_red/layer_30/layer_30_heatmap.png}
    \caption{layer 30：CLS：中心コサイン類似度ヒートマップ}
  \end{subfigure}\hfill
  \begin{subfigure}[t]{0.48\linewidth}
    \centering
    \includegraphics[width=\linewidth]{figures/CCS_heatmaps_cls_thesis_red/layer_final/layer_final_heatmap.png}
    \caption{layer final：CLS：中心コサイン類似度ヒートマップ}
  \end{subfigure}

  \vspace{0.6em}

  % 2行目
  \begin{subfigure}[t]{0.48\linewidth}
    \centering
    \includegraphics[width=\linewidth]{figures/umap_cls_thesis/layer_30/umap_cls_by_aa.png}
    \caption{layer 30：CLS：UMAP可視化（AA）}
  \end{subfigure}\hfill
  \begin{subfigure}[t]{0.48\linewidth}
    \centering
    \includegraphics[width=\linewidth]{figures/umap_cls_thesis/layer_final/umap_cls_by_aa.png}
    \caption{layer final：CLS：UMAP可視化（AA）}
  \end{subfigure}

  \vspace{0.6em}

  % 3行目
  \begin{subfigure}[t]{0.48\linewidth}
    \centering
    \includegraphics[width=\linewidth]{figures/umap_cls_thesis/layer_30/umap_cls_by_domain.png}
    \caption{layer 30：CLS：UMAP可視化（Domain）}
  \end{subfigure}\hfill
  \begin{subfigure}[t]{0.48\linewidth}
    \centering
    \includegraphics[width=\linewidth]{figures/umap_cls_thesis/layer_final/umap_cls_by_domain.png}
    \caption{layer final：CLS：UMAP可視化（Domain）}
  \end{subfigure}

  \caption{layer 30とlayer finalのCLSトークンに関する中心コサイン類似度ヒートマップおよびUMAP可視化}
  \label{fig:cls30_final}
\end{figure}

\subsubsection{各塩基トークンにおけるUMAP可視化}

\begin{figure}[H]
  \centering

  % 1行目
  \begin{subfigure}[t]{0.48\linewidth}
    \centering
    \includegraphics[width=\linewidth]{figures/umap_token_thesis_BASE/layer_00/umap_token_by_base.png}
    \caption{layer 0：各塩基トークン：UMAP可視化(Base)}
  \end{subfigure}\hfill
  \begin{subfigure}[t]{0.48\linewidth}
    \centering
    \includegraphics[width=\linewidth]{figures/umap_token_thesis_BASE/layer_05/umap_token_by_base.png}
    \caption{layer 5：各塩基トークン：UMAP可視化(Base)}
  \end{subfigure}

  \vspace{0.6em}

  % 2行目
  \begin{subfigure}[t]{0.48\linewidth}
    \centering
    \includegraphics[width=\linewidth]{figures/umap_token_thesis/layer_00/umap_token_by_domain.png}
    \caption{layer 0：各塩基トークン：UMAP可視化（Domain）}
  \end{subfigure}\hfill
  \begin{subfigure}[t]{0.48\linewidth}
    \centering
    \includegraphics[width=\linewidth]{figures/umap_token_thesis/layer_05/umap_token_by_domain.png}
    \caption{layer 5：各塩基トークン：UMAP可視化（Domain）}
  \end{subfigure}

  \vspace{0.6em}

  % 3行目
  \begin{subfigure}[t]{0.48\linewidth}
    \centering
    \includegraphics[width=\linewidth]{figures/umap_token_thesis/layer_00/umap_token_by_structure.png}
    \caption{layer 0：各塩基トークン：UMAP可視化（structure）}
  \end{subfigure}\hfill
  \begin{subfigure}[t]{0.48\linewidth}
    \centering
    \includegraphics[width=\linewidth]{figures/umap_token_thesis/layer_05/umap_token_by_structure.png}
    \caption{layer 5：各塩基トークン：UMAP可視化（structure）}
  \end{subfigure}

  \caption{layer0とlayer5の各塩基トークンに関するUMAP可視化}
  \label{fig:tokens0_5}
\end{figure}


\begin{figure}[H]
  \centering

  % 1行目
  \begin{subfigure}[t]{0.48\linewidth}
    \centering
    \includegraphics[width=\linewidth]{figures/umap_token_thesis_BASE/layer_10/umap_token_by_base.png}
    \caption{layer 10：各塩基トークン：UMAP可視化（Base）}
  \end{subfigure}\hfill
  \begin{subfigure}[t]{0.48\linewidth}
    \centering
    \includegraphics[width=\linewidth]{figures/umap_token_thesis_BASE/layer_15/umap_token_by_base.png}
    \caption{layer 15：各塩基トークン：UMAP可視化（Base）}
  \end{subfigure}

  \vspace{0.6em}

  % 2行目
  \begin{subfigure}[t]{0.48\linewidth}
    \centering
    \includegraphics[width=\linewidth]{figures/umap_token_thesis/layer_10/umap_token_by_domain.png}
    \caption{layer 10：各塩基トークン：UMAP可視化（Domain）}
  \end{subfigure}\hfill
  \begin{subfigure}[t]{0.48\linewidth}
    \centering
    \includegraphics[width=\linewidth]{figures/umap_token_thesis/layer_15/umap_token_by_domain.png}
    \caption{layer 15：各塩基トークン：UMAP可視化（Domain）}
  \end{subfigure}

  \vspace{0.6em}

  % 3行目
  \begin{subfigure}[t]{0.48\linewidth}
    \centering
    \includegraphics[width=\linewidth]{figures/umap_token_thesis/layer_10/umap_token_by_structure.png}
    \caption{layer 10：各塩基トークン：UMAP可視化（structure）}
  \end{subfigure}\hfill
  \begin{subfigure}[t]{0.48\linewidth}
    \centering
    \includegraphics[width=\linewidth]{figures/umap_token_thesis/layer_15/umap_token_by_structure.png}
    \caption{layer 15：各塩基トークン：UMAP可視化（structure）}
  \end{subfigure}

  \caption{layer10とlayer15の各塩基トークンに関するUMAP可視化}
  \label{fig:tokens10_15}
\end{figure}

\begin{figure}[H]
  \centering

  % 1行目
  \begin{subfigure}[t]{0.48\linewidth}
    \centering
    \includegraphics[width=\linewidth]{figures/umap_token_thesis_BASE/layer_20/umap_token_by_base.png}
    \caption{layer 20：各塩基トークン：UMAP可視化（Base）}
  \end{subfigure}\hfill
  \begin{subfigure}[t]{0.48\linewidth}
    \centering
    \includegraphics[width=\linewidth]{figures/umap_token_thesis/layer_22/umap_token_by_aa.png}
    \caption{layer 22：各塩基トークン：UMAP可視化（AA）}
  \end{subfigure}

  \vspace{0.6em}

  % 2行目
  \begin{subfigure}[t]{0.48\linewidth}
    \centering
    \includegraphics[width=\linewidth]{figures/umap_token_thesis/layer_20/umap_token_by_domain.png}
    \caption{layer 20：各塩基トークン：UMAP可視化（Domain）}
  \end{subfigure}\hfill
  \begin{subfigure}[t]{0.48\linewidth}
    \centering
    \includegraphics[width=\linewidth]{figures/umap_token_thesis/layer_22/umap_token_by_domain.png}
    \caption{layer 22：各塩基トークン：UMAP可視化（Domain）}
  \end{subfigure}

  \vspace{0.6em}

  % 3行目
  \begin{subfigure}[t]{0.48\linewidth}
    \centering
    \includegraphics[width=\linewidth]{figures/umap_token_thesis/layer_20/umap_token_by_structure.png}
    \caption{layer 20：各塩基トークン：UMAP可視化（structure）}
  \end{subfigure}\hfill
  \begin{subfigure}[t]{0.48\linewidth}
    \centering
    \includegraphics[width=\linewidth]{figures/umap_token_thesis/layer_22/umap_token_by_structure.png}
    \caption{layer 22：各塩基トークン：UMAP可視化（structure）}
  \end{subfigure}

  \caption{layer20とlayer22の各塩基トークンに関するUMAP可視化}
  \label{fig:tokens20_22}
\end{figure}

\begin{figure}[H]
  \centering

  % 1行目
  \begin{subfigure}[t]{0.48\linewidth}
    \centering
    \includegraphics[width=\linewidth]{figures/umap_token_thesis/layer_23/umap_token_by_aa.png}
    \caption{layer 23：各塩基トークン：UMAP可視化（AA）}
  \end{subfigure}\hfill
  \begin{subfigure}[t]{0.48\linewidth}
    \centering
    \includegraphics[width=\linewidth]{figures/umap_token_thesis/layer_24/umap_token_by_aa.png}
    \caption{layer 24：各塩基トークン：UMAP可視化（AA）}
  \end{subfigure}

  \vspace{0.6em}

  % 2行目
  \begin{subfigure}[t]{0.48\linewidth}
    \centering
    \includegraphics[width=\linewidth]{figures/umap_token_thesis/layer_23/umap_token_by_domain.png}
    \caption{layer 23：各塩基トークン：UMAP可視化（Domain）}
  \end{subfigure}\hfill
  \begin{subfigure}[t]{0.48\linewidth}
    \centering
    \includegraphics[width=\linewidth]{figures/umap_token_thesis/layer_24/umap_token_by_domain.png}
    \caption{layer 24：各塩基トークン：UMAP可視化（Domain）}
  \end{subfigure}

  \vspace{0.6em}

  % 3行目
  \begin{subfigure}[t]{0.48\linewidth}
    \centering
    \includegraphics[width=\linewidth]{figures/umap_token_thesis/layer_23/umap_token_by_structure.png}
    \caption{layer 23：各塩基トークン：UMAP可視化（structure）}
  \end{subfigure}\hfill
  \begin{subfigure}[t]{0.48\linewidth}
    \centering
    \includegraphics[width=\linewidth]{figures/umap_token_thesis/layer_24/umap_token_by_structure.png}
    \caption{layer 24：各塩基トークン：UMAP可視化（structure）}
  \end{subfigure}

  \caption{layer23とlayer24の各塩基トークンに関するUMAP可視化}
  \label{fig:tokens23_24}
\end{figure}


\begin{figure}[H]
  \centering

  % 1行目
  \begin{subfigure}[t]{0.48\linewidth}
    \centering
    \includegraphics[width=\linewidth]{figures/umap_token_thesis/layer_26/umap_token_by_aa.png}
    \caption{layer 26：各塩基トークン：UMAP可視化（AA）}
  \end{subfigure}\hfill
  \begin{subfigure}[t]{0.48\linewidth}
    \centering
    \includegraphics[width=\linewidth]{figures/umap_token_thesis/layer_28/umap_token_by_aa.png}
    \caption{layer 28：各塩基トークン：UMAP可視化（AA）}
  \end{subfigure}

  \vspace{0.6em}

  % 2行目
  \begin{subfigure}[t]{0.48\linewidth}
    \centering
    \includegraphics[width=\linewidth]{figures/umap_token_thesis/layer_26/umap_token_by_domain.png}
    \caption{layer 26：各塩基トークン：UMAP可視化（Domain）}
  \end{subfigure}\hfill
  \begin{subfigure}[t]{0.48\linewidth}
    \centering
    \includegraphics[width=\linewidth]{figures/umap_token_thesis/layer_28/umap_token_by_domain.png}
    \caption{layer 28：各塩基トークン：UMAP可視化（Domain）}
  \end{subfigure}

  \vspace{0.6em}

  % 3行目
  \begin{subfigure}[t]{0.48\linewidth}
    \centering
    \includegraphics[width=\linewidth]{figures/umap_token_thesis/layer_26/umap_token_by_structure.png}
    \caption{layer 26：各塩基トークン：UMAP可視化（structure）}
  \end{subfigure}\hfill
  \begin{subfigure}[t]{0.48\linewidth}
    \centering
    \includegraphics[width=\linewidth]{figures/umap_token_thesis/layer_28/umap_token_by_structure.png}
    \caption{layer 28：各塩基トークン：UMAP可視化（structure）}
  \end{subfigure}

  \caption{layer26とlayer28の各塩基トークンに関するUMAP可視化}
  \label{fig:tokens26_28}
\end{figure}

\begin{figure}[H]
  \centering

  % 1行目
  \begin{subfigure}[t]{0.48\linewidth}
    \centering
    \includegraphics[width=\linewidth]{figures/umap_token_thesis/layer_30/umap_token_by_aa.png}
    \caption{layer 30：各塩基トークン：UMAP可視化（AA）}
  \end{subfigure}\hfill
  \begin{subfigure}[t]{0.48\linewidth}
    \centering
    \includegraphics[width=\linewidth]{figures/umap_token_thesis/layer_final/umap_token_by_aa.png}
    \caption{layer final：各塩基トークン：UMAP可視化（AA）}
  \end{subfigure}

  \vspace{0.6em}

  % 2行目
  \begin{subfigure}[t]{0.48\linewidth}
    \centering
    \includegraphics[width=\linewidth]{figures/umap_token_thesis/layer_30/umap_token_by_domain.png}
    \caption{layer 30：各塩基トークン：UMAP可視化（Domain）}
  \end{subfigure}\hfill
  \begin{subfigure}[t]{0.48\linewidth}
    \centering
    \includegraphics[width=\linewidth]{figures/umap_token_thesis/layer_final/umap_token_by_domain.png}
    \caption{layer final：各塩基トークン：UMAP可視化（Domain）}
  \end{subfigure}

  \vspace{0.6em}

  % 3行目
  \begin{subfigure}[t]{0.48\linewidth}
    \centering
    \includegraphics[width=\linewidth]{figures/umap_token_thesis/layer_30/umap_token_by_structure.png}
    \caption{layer 30：各塩基トークン：UMAP可視化（structure）}
  \end{subfigure}\hfill
  \begin{subfigure}[t]{0.48\linewidth}
    \centering
    \includegraphics[width=\linewidth]{figures/umap_token_thesis/layer_final/umap_token_by_structure.png}
    \caption{layer final：各塩基トークン：UMAP可視化（structure）}
  \end{subfigure}

  \caption{layer 30とlayer finalの各塩基トークンに関するUMAP可視化}
  \label{fig:tokens30_final}
\end{figure}
\subsection{AttnLRPによる解析}

\subsubsection{層ごとのトークンタイプ別寄与率}

\begin{figure}[H]
  \centering
  \includegraphics[width=0.85\linewidth]{figures/AttnLRP/thesis_absolute.png}
  \caption{層ごとのトークンタイプ別寄与率(絶対値)}
  \label{fig:attnlrp_token_absolute}
\end{figure}

\begin{figure}[H]
  \centering
  \includegraphics[width=0.85\linewidth]{figures/AttnLRP/thesis_ratio.png}
  \caption{層ごとのトークンタイプ別寄与率(比率)}
  \label{fig:attnlrp_tokentype_ratio}
\end{figure}

層を遡るにつれて（出力層から入力層へ向かうにつれて）Relevanceの総和は減少していることがわかる。
また層が進むにつれて、CLSのrelevanceの割合が単調に増えている。
どの層においてもEOSのrelevanceは非常に小さいことがわかる。

\subsubsection{層ごとのトークンタイプ別コンポーネント割合}

% --- CLS: 絶対値 ---
\begin{figure}[H]
  \centering
  \includegraphics[width=0.85\linewidth]{figures/AttnLRP/AttnLRP_cls_absolute.png}
  \caption{層ごとのCLSのコンポーネント寄与率(絶対値)}
  \label{fig:attnlrp_cls_absolute}
\end{figure}

% --- CLS: 比率 ---
\begin{figure}[H]
  \centering
  \includegraphics[width=0.85\linewidth]{figures/AttnLRP/AttnLRP_cls_ratio.png}
  \caption{層ごとのCLSのコンポーネント寄与率(比率)}
  \label{fig:attnlrp_cls_ratio}
\end{figure}

層が進むにつれて、CLSのrelevanceの割合が単調に増えていることがわかる。
またその中でもResidual成分が多数の層において大きな割合を占めていることがわかる。

% --- Tokens: 絶対値 ---
\begin{figure}[H]
  \centering
  \includegraphics[width=0.85\linewidth]{figures/AttnLRP/AttnLRP_tokens_absolute.png}
  \caption{層ごとの各塩基トークンのコンポーネント割合(絶対値)}
  \label{fig:attnlrp_tokens_absolute}
\end{figure}

% --- Tokens: 比率 ---
\begin{figure}[H]
  \centering
  \includegraphics[width=0.85\linewidth]{figures/AttnLRP/AttnLRP_tokens_ratio.png}
  \caption{層ごとの各塩基トークンのコンポーネント割合(比率)}
  \label{fig:attnlrp_tokens_ratio}
\end{figure}

各塩基トークンのrelevanceの総和は層が進むにつれて、増加していくと考えられたが、24層を境に減少していることがわかる。
CLSと同様に多数の層においてResidual成分が大きな割合を占めていることがわかる。


\subsubsection{ラベル決定因子について}

% GGGCTCGTAGATCAGTGGTAGATCGACACCTTCGCAAGGTGTAGGCCCCGGGTTCAAATCCCGGCGAGTCCA
\begin{figure}[H]
  \centering
  \includegraphics[width=0.85\linewidth]{figures/arc_Ala/1.png}
  \caption{arc\_Alaのrelevanceフローグラフ}
  \label{fig:attnlrp_arc_Ala}
\end{figure}

\begin{table}[H]
  \centering
  \begin{tabular}{c c r r r r r}
    \hline
    Token & Pos & Diff (Z) & p-value & Late Avg(Z) & Early Avg(Z) \\
    \hline
    C     & 36  & 2.0377   & 0.0027  & 4.5758      & 2.5382       \\
    A     & 37  & 1.9308   & 0.0027  & 2.3809      & 0.4501       \\
    T     & 70  & 1.1949   & 9.5e-06 & 0.2730      & -0.9219      \\
    A     & 57  & 1.1695   & 0.0029  & 0.9296      & -0.2399      \\
    A     & 44  & 1.0999   & 2.5e-05 & 0.2755      & -0.8244      \\
    T     & 5   & 0.9975   & 1.9e-07 & 0.4413      & -0.5561      \\
    G     & 10  & 0.7326   & 0.0027  & 0.5168      & -0.2158      \\
    C     & 47  & 0.6545   & 1.4e-05 & 0.1185      & -0.5360      \\
    \hline
  \end{tabular}
  \caption{arc\_Alaの層間比較検定}
  \label{tab:token_arc_stats}
\end{table}



\begin{figure}[H]
  \centering
  \includegraphics[width=0.85\linewidth]{figures/bac_Ala/1.png}
  \caption{bac\_Alaのrelevanceフローグラフ}
  \label{fig:attnlrp_bac_Ala}
\end{figure}

% GGGCCCGTAGTTTACCGGTAAAATATCTGGCTGGCAGTCAGAAGTAAGGGGTTCAATTCCCCTCGGGTCCACCA

\begin{table}[H]
  \centering
  \begin{tabular}{c c r r r r r}
    \hline
    Token & Pos & Diff (Z) & p-value & Late Avg(Z) & Early Avg(Z) \\
    \hline
    C     & 36  & 1.4633   & 0.0055  & 3.8786      & 2.4153       \\
    T     & 13  & 0.7319   & 9.5e-05 & 0.1394      & -0.5924      \\
    T     & 70  & 0.6270   & 0.0125  & 0.2154      & -0.4116      \\
    C     & 4   & 0.6034   & 2.5e-04 & 0.5368      & -0.0666      \\
    C     & 6   & 0.5775   & 1.6e-06 & 0.1435      & -0.4340      \\
    T     & 12  & 0.5204   & 1.3e-04 & 0.1247      & -0.3957      \\
    G     & 10  & 0.4741   & 0.0020  & 0.2780      & -0.1961      \\
    \hline
  \end{tabular}
  \caption{bac\_Alaの層間比較検定}
  \label{tab:token_bac_stats}
\end{table}

% GGGCGCTTGGTCTAGGGGTATGATCTTCGTTTCGCACCACACCGCGTCTACGAAAAGTCGTGGGTTCGATTCCTACAGCGTCCA
\begin{figure}[H]
  \centering
  \includegraphics[width=0.85\linewidth]{figures/euk_Ala/1.png}
  \caption{euk\_Alaのrelevanceフローグラフ}
  \label{fig:attnlrp_euk_Ala}
\end{figure}

\begin{table}[H]
  \centering
  \begin{tabular}{c c r r r r}
    \hline
    Token & Pos & Diff (Z) & p-value & Late Avg(Z) & Early Avg(Z) \\
    \hline
    G     & 3   & 1.2647   & 4.8e-06 & 1.3605      & 0.0958       \\
    A     & 37  & 1.2384   & 0.0281  & 2.0985      & 0.8601       \\
    T     & 32  & 1.1509   & 0.0073  & 1.7093      & 0.5584       \\
    T     & 31  & 0.8874   & 7.4e-04 & 1.1694      & 0.2820       \\
    G     & 9   & 0.6027   & 1.8e-04 & 0.6402      & 0.0375       \\
    T     & 11  & 0.5637   & 3.3e-05 & 0.6097      & 0.0461       \\
    G     & 5   & 0.5188   & 5.9e-04 & 0.4383      & -0.0805      \\
    \hline
  \end{tabular}
  \caption{euk\_Alaの層間比較検定}
  \label{tab:token_euk_stats}
\end{table}

\clearpage

\section{考察}

\subsection{分類モデル}
\subsubsection{モデルの全体指標}

結果で示した通り、モデル全体として優れた分類性能を示していることが分かる。
またTop-3 Accuracyも99.47\%に達しており、これはアミノ酸に関しては予測できているがdomainの識別が難しいケースが多数存在するのではないかと示唆される。
これについては\ref{Consideration-top10-misclassifications}で詳しく考察する。

\subsubsection{ラベルごとの評価}
\label{per_label}
表\ref{tab:rinalmo_test_per_label}より、support数が多いラベルに関しては比較的高い性能を示しているが、support数が少ないラベルに関しては性能が著しく低下していることが分かる。
これに関しては以下の二つの原因が考えられる。

\textbf{(1)最適化の影響}\\
本分類モデルではでは、式\eqref{eq:cross-entropy-loss}で表される交差エントロピー損失を最小化するようにモデルの学習は行われる。
ここで \(N\) は全サンプル数であり、各サンプルの損失が同じ重みで平均化される。
したがって、学習データにおいてラベル \(y\) の出現比率 \(\pi_y=\frac{N_y}{N}\)（\(N_y\): ラベル \(y\) のサンプル数）が偏っている場合、損失の期待値は
\begin{equation}
  \mathbb{E}[L] \approx \sum_{y}\pi_y\,\mathbb{E}_{x\sim p(x\mid y)}\left[-\log p_\theta(y\mid x)\right]
  \label{eq:ce_decompose}
\end{equation}
のように，各クラスの寄与が \(\pi_y\) によって重み付けされた形で現れる。このため、多数派ラベル（\(\pi_y\) が大きい）の誤分類をわずかに減らすだけでも全体損失が大きく下がりやすく、最適化の過程で多数派を優先する更新が起こりやすい。
極端な場合、多数派のみを予測しても見かけの正解率が高くなり得るため、学習が進んでいるように見えても少数派の再現率（Recall）が低い状態が残りやすい。

また，予測確率 \(p_\theta(y\mid x)\) は学習データの分布を反映するため，事前確率 \(p(y)\) が小さい少数派ラベルは，入力が同程度に曖昧な領域では選ばれにくい．ベイズ則
\begin{equation}
  p(y\mid x)=\frac{p(x\mid y)p(y)}{p(x)}
  \label{eq:bayes}
\end{equation}
より，クラス条件付き分布 \(p(x\mid y)\) が同程度でも \(p(y)\) が小さいほど \(p(y\mid x)\) は小さくなりやすい。
この性質は，本モデルが直接 \(p_\theta(y\mid x)\) を学習する場合でも、データ比率を通じて暗黙に反映され、少数ラベルが不利になりやすいと考えられる。

\medskip
\textbf{(2)ラベルごとの識別要素の欠如}\\
(1)が主に「最適化の影響」であるのに対し、(2)は「少数派クラスの情報不足による推定の難しさ」に対応する。
少数派ラベルは \(N_y\) が小さいため、そのクラスに固有な特徴を十分に観測できず、クラス内の多様性が学習データに現れにくい。
その結果、モデルが学習する表現は多数派の特徴に引っ張られやすく、少数派を区別するための安定した識別要素が形成されにくい。

統計的には、少数派クラスに関する推定はサンプル数が少ないほど不確実性が大きくなり、汎化誤差が増大しやすい。
直感的には、少数派の「典型例」や「例外」が十分に見えないため、学習した決定境界が少数派領域を適切に覆えず、多数派側へ押し込まれる形になりやすい。


したがって、ラベル不均衡データによる学習では(1)の最適化上の偏りに加え、(2)のデータ不足による識別要素の欠如が重なって、少数派ラベルほど予測が難しくなると考察することができる。

また少数ラベルに関しては、Precisionが高く、Recallが低くなることが多くなると考えられる。これは表\ref{tab:ovr_confusion},式\eqref{eq:precison},式\eqref{eq:recall}から以下のように考えることができる。
これは少数ラベルは$\mathrm{FN}_c$が多くなり、Recallが低くなる一方で、$\mathrm{FP}_c$が少ないためPrecisionが高くなるためである。

ただし、arc\_Trp(support=19),euk\_iMet(support=21)などsupport数が少ないラベルでも高い性能を示している例も存在している。これらのラベルに関しては、配列に他のラベルと明確に異なる特徴が存在しており、少数サンプルでも識別要素を学習しやすかった可能性が考えられる。

support数から考えて、相対的に難易度が高いクラスとしては、euk\_Met（Recall=0.6706, F1=0.7500, Support=85）が挙げられる。Recallが低く、正例を取りこぼす傾向が示唆される。これらのクラスは他クラスと配列特徴が近接しやすい可能性があり、モデルが保守的になることでRecallが低下していることが考えられる。



\subsubsection{誤分類Top10}
\label{Consideration-top10-misclassifications}

同一アミノ酸を維持したままドメインのみが置換されるパターンが大半を占めていて、これはモデルがアンチコドンなど、アミノ酸に重要な配列特徴を捉えていることを示唆しているが、一方でドメイン識別に関わる特徴の学習は相対的に難しいことも示唆される。

真核由来クラス（euk\_*）が細菌由来クラス（bac\_*）へ誤分類される方向が上位を占める点は重要である。
これは多数ラベルである細菌クラスに引っ張られる最適化上の影響（\ref{per_label}(1)）が関与している可能性がある。
また、真核と細菌のtRNA配列は進化的に近い関係にあり、配列特徴が類似している場合も多いため、モデルが両者を区別するための安定した識別要素を学習しにくい（\ref{per_label}(2)）ことも考えられる。

一方で、bac\_Phe$\rightarrow$euk\_Phe（18件）やbac\_Leu$\rightarrow$euk\_Leu（49件）のように、eukaryotaクラスがbacteriaクラスへ誤分類されるパターンも確認された。
従って、誤分類は真核$\rightarrow$細菌に限定されず、少なくとも一部のアミノ酸については、ドメイン差よりもアミノ酸特異的特徴の方が支配的に学習されている。

従来、aaRSやtRNA 遺伝子は翻訳系の中核に位置し、細胞内の多くの要素と整合的に機能する必要があることから、基本的には垂直伝播に従うと考えられてきた。
しかし近年、これらの遺伝子群においても水平伝播が起こり得ることが報告されており\cite{woese2000aminoacyl}\cite{tuller2011association}、翻訳系であっても進化史が必ずしも純粋な系統樹として記述できないことが示されている。
この点を踏まえると、本研究で観察された少数の誤分類は、単なるモデル性能の限界というよりも、水平伝播の影響が反映されている可能性がある。

例えば、tRNA が水平伝播により外来由来の配列としてゲノムに導入されると、受け手側が本来もつドメイン固有の配列傾向とは異なる「外来の配列パターン」が混入し得る。
本分類モデルは、後述する\ref{subsec:internal_analysis} で考察するように、学習過程でドメインに特徴的なモチーフ分布などを手がかりとして分類していると考えられるため、外来由来のtRNAは実際の所属ドメインとは異なるドメインに近い特徴を示し、誤分類として現れることがあり得る。
さらに、aaRS と tRNA は認識関係を介して共進化しているため、aaRS 側の水平伝播や置換が生じた場合、対応する tRNA 側でも同化や部分的な配列改変が進み、結果として配列特徴が「ドナー由来の傾向」と「受け手側での適応」の混合状態となる可能性がある。

またtRNAにおいては遺伝子変換が起こることも先行研究\cite{amstutz1985concerted}などから知られている。

このようにtRNAに関しては、水平伝播や遺伝子変換など、垂直伝播の枠組みを超えた進化過程が存在する可能性が示されている。

以上より、本研究の誤分類例は「モデルが誤っている」ことを示すというよりも、モデルが配列に含まれる系統的起源に関する情報を学習しているのではないかと考えられる。
従って、誤分類サンプルを個別に検討することは、水平伝播や遺伝子置換が疑われる候補を抽出し、翻訳系遺伝子の例外的な進化過程を検討するための手がかりになり得る。


以上を総合すると、本モデルの誤りは「アミノ酸を取り違える」よりも「同一アミノ酸のままドメインを取り違える」形で多発しており、本モデルがtRNAのアミノ酸同定に関わる主要な配列特徴を高精度に捉えていることを裏付ける結果といえる。
一方で、ドメイン識別に寄与する特徴の抽出は相対的に難しく、特にeukaryotaクラスがbacteriaクラスへ吸収される傾向が示唆された。
また誤分類はモデルの性質だけに留まらず、aaRS,tRNAの水平伝播に関係する可能性も示唆される。

\subsection{内部表現解析}
\label{subsec:internal_analysis}

本節では、モデルが層を重ねるにつれて配列表現がどのように再編成されるかを調べるため、
(1) CLSトークン埋め込み、(2) 各塩基トークン埋め込みを対象に解析した。
CLSについてはUMAP可視化に加えて中心化コサイン類似度を用い、各塩基トークンについてはUMAP可視化に基づいて層ごとの変化を比較した。

\subsubsection{CLSトークン}
\label{subsec:cls_analysis}

CLS埋め込みの層別解析から、表現は
\textbf{(i) ラベル非依存な汎用表現} $\rightarrow$
\textbf{(ii) ドメイン情報の整理} $\rightarrow$
\textbf{(iii) ドメインを超えたアミノ酸中心表現}
へと段階的に移行することが示された。

\paragraph{浅層（0--10層）：汎用的表現段階}

0層(図\ref{fig:cls0_5}(a))では中心化コサイン類似度が全体に高く、サンプル間の表現が広く似通っていた。
UMAP(図\ref{fig:cls0_5}(c),(e))でもドメイン別・アミノ酸別のいずれでも明確な分離は見られず、CLSがまだラベルに対応した構造を獲得していない段階と解釈できる。
5層(図\ref{fig:cls0_5}(b),(d),(f))・10層(図\ref{fig:cls10_15}(a),(c),(e))でも傾向は概ね同様であり、ドメイン分離がわずかに進む兆候はあるものの、識別的情報は十分に顕在化していない。

\paragraph{中層前半（15層）：表現空間の転換点}

15層ではUMAP上(図\ref{fig:cls10_15}(f))でドメイン別クラスタが比較的明瞭になり、CLSがドメイン情報を反映し始めた。
また、Archaea由来の一部ではアミノ酸別の分離も見られ、アミノ酸識別に関わる特徴が部分的に立ち上がり始めた可能性がある。
一方、中心化コサイン類似度(図\ref{fig:cls10_15}(b))ではドメイン間差はまだ明確ではなく、分離は移行段階にあると考えられる。

\paragraph{中層後半（20--22層）：ドメイン表現の確立}

20層ではUMAP(図\ref{fig:cls20_22}(a),(c),(e))でドメイン分離が明確となり、CLS表現がドメイン情報を強く担う段階に到達した。
中心化コサイン類似度でもドメイン内類似度が上昇し、ドメイン間との差が現れ始めた(図\ref{fig:cls20_22}(a)の赤枠で囲った対角ブロックはドメインごとのコサイン類似度である)。
特にArchaeaではアミノ酸別分離も顕著であり、アミノ酸識別特徴の形成が先行して進む可能性が示唆された。
22層(図\ref{fig:cls20_22}(b),(d),(f))ではこの傾向がさらに安定し、bacteria/eukaryotaでもアミノ酸別の小規模クラスタが徐々に観察され、次段階への準備が進んでいることが示された。

\paragraph{深層（23--final層）：アミノ酸中心表現への移行}

23層(図\ref{fig:cls23_24}(a),(c),(e))以降、bacteria/eukaryotaでもアミノ酸別クラスタが急速に明確化した。
中心化コサイン類似度でも、ドメインをまたいだ同一アミノ酸間の類似度が相対的に高まり、
表現の主軸がドメインからアミノ酸へ移行し始めたことが示唆された(表\ref{fig:cls23_24}の赤枠のブロックはアミノ酸が同じラベルの類似度を指している)。
26層(図\ref{fig:cls26_28}(a),(c),(e))ではこの傾向がより強まり、ドメイン内類似度は低下する一方でアミノ酸別のまとまりが強くなった。
28層(図\ref{fig:cls26_28}(b),(d),(f))・30層(図\ref{fig:cls30_final}(a),(c),(e))・final層(図\ref{fig:cls30_final}(b),(d),(f))では大きな構造は維持されつつ、ドメイン寄与がさらに弱まり、アミノ酸間の連続的な類似構造（クラスタ間距離の接近）が強化された。

\paragraph{まとめ}
以上より、本モデルのCLS表現は
中層でドメイン情報を整理・確立した後、
深層でドメインを超えたアミノ酸中心表現へと再編成され、ラベル分類に最適化されるという段階的な遷移を示した。

\subsubsection{各塩基トークン}
\label{subsec:token_analysis}

各層の各塩基トークン埋め込みをUMAPで可視化した結果、トークン表現は主に
\textbf{structure}を軸として分離が進み、深層でより細かな部位差へと分解される傾向が観察された。

\paragraph{浅層（0--15層）：分離の未発達}
0層(図\ref{fig:tokens0_5}(a),(c),(e))では、ドメイン、structure、塩基の明確な分離は見られず、
トークン表現はラベル対応構造をまだ獲得していない段階と考えられる。

5層(図\ref{fig:tokens0_5}(b),(d),(f))、10層(図\ref{fig:tokens10_15}(a),(c),(e))では塩基ごとに明確な分離が見られたが、structure別・ドメイン別の分離はまだ不明瞭であった。
15層(図\ref{fig:tokens10_15}(b),(d),(f))ではクラスタ分離が進み、塩基を軸として、structureごとの分離が確認された。

\paragraph{中層（20--23層）：structure主導の分離とドメインの付随}
20層(図\ref{fig:tokens20_22}(a),(c),(e))ではstructure別クラスタ分離が明確となり、その内部でdomainごとのまとまりもある程度観察された。
15層の時点ではstructure内で塩基ごとに分かれていたが、20層では塩基ごとのまとまりが薄れ、structureを軸にした分離が優勢になっていることがわかる。
これは、まずstructureに関する特徴が前景化し、その上でdomain差が付随的に現れていることを示唆する。
22層(図\ref{fig:tokens20_22}(b),(d),(f))、23層(図\ref{fig:tokens23_24}(a),(c),(e))では塩基ごとの分離がさらに弱まり、structure内でdomain別の小クラスタがより明瞭に観察された。

\paragraph{深層（24--final層）：細粒度な部位差への分解}
24層(図\ref{fig:tokens23_24}(b),(d),(f))ではクラスタがさらに細分化し、特にanticodon\_coreが小クラスタとして分離する傾向が見られた。
26層(図\ref{fig:tokens26_28}(a),(c),(e))では細分化がより進み、anticodon\_coreに加えてacceptor\_stem\_5side、anticodon\_loop\_5sideなども
より独立した小クラスタとして観察された。
28層(図\ref{fig:tokens26_28}(b),(d),(f))は26層から大きな変化は少ないが、30層(図\ref{fig:tokens30_final}(b),(d),(f))ではT\_loopがより小さなクラスタとして分離する傾向が見られた。

\paragraph{まとめ}
トークン表現はまずは塩基ごとに分かれ、その後はstructureを軸にした分離が進み、その中でもdomainによって小クラスタが形成されることが分かった。
深層ではstructure内の細粒度な部位差へと段階的に分解される傾向を示した。

\subsubsection{内部表現解析のまとめ}
\label{sec:transision_analysis}

\ref{subsec:cls_analysis},\ref{subsec:token_analysis}では、CLS、配列トークンの双方において、層における学習推移が確認された。
このような層の学習推移に関しては、BERT型のモデルが層ごとの学習推移を持つことを示した先行研究\cite{jawahar2019does}と整合的である。
一方でなぜこのような学習推移が生じるのかを完全に説明する枠組みは未だ確立されておらず本研究でも結果に整合する形で限定的に解釈を試みる。

層の学習推移に関して2つの観点から考察を行う。

\paragraph{(1)モデルの構造}

分類に用いるベクトルは最終的に\ref{sec:classification-model}より $\mathbf{h}\in\mathbb{R}^{d}$（$\mathbf{h}=\mathbf{h}_{\mathrm{CLS}}$, $d=256$）と表せる。
logitsは以下の様に表すことができる。
\begin{equation}
  z_k = \mathbf{w}_k^\top \mathbf{h},\qquad k\in\{1,\dots,K\},\ \ K=68
\end{equation}
を得る。ここで $\mathbf{w}_k\in\mathbb{R}^{d}$ はクラス $k$ に対応する重みベクトルである。
交差エントロピー損失は式\eqref{eq:cross-entropy-loss}より与えられ、以下のように変形することができる。
\begin{equation}
  \mathcal{L}
  = -z_{y_i} + \log\sum_{k=1}^{K}e^{z_k},
  \label{eq:ce_ja}
\end{equation}
ただし $\mathbf{z}=(z_1,\dots,z_K)$ とする。

softmaxを用いて確率を
\begin{equation}
  p_k=\mathrm{softmax}(\mathbf{z})_k=\frac{e^{z_k}}{\sum_{i=1}^{K}e^{z_i}}
\end{equation}
と定義する。このときロジットに関する勾配は既知の形として
\begin{equation}
  \frac{\partial \mathcal{L}}{\partial z_k} = p_k - \mathbb{1}[k=y]
  \label{eq:dz_ja}
\end{equation}
となる。モデルの学習初期では $p_k\approx 1/K$ が成り立つと近似できるため
\begin{equation}
  \frac{\partial \mathcal{L}}{\partial z_k}\approx
  \begin{cases}
    \frac{1}{K}-1=-\frac{K-1}{K} & (k=y),    \\[4pt]
    \frac{1}{K}                  & (k\neq y)
  \end{cases}
  \label{eq:dz_init_ja}
\end{equation}
を得る。特に $K=68$ では $\partial \mathcal{L}/\partial z_y\approx -67/68\approx -0.985$
$\partial \mathcal{L}/\partial z_{k\neq y}\approx 1/68\approx 0.015$ である。
負の大きな誤差は正解クラス1つに集中する一方、正の誤差は残り $K-1$ 個の非正解クラスへ小さく分散する。

$z_k=\mathbf{w}_k^\top\mathbf{h}$ より連鎖律を用いると，
\begin{align}
  \frac{\partial \mathcal{L}}{\partial \mathbf{w}_k}
   & =\frac{\partial \mathcal{L}}{\partial z_k}\frac{\partial z_k}{\partial \mathbf{w}_k}
  =(p_k-\mathbb{1}[k=y])\,\mathbf{h},
  \label{eq:dw_ja}                                                                                    \\
  \frac{\partial \mathcal{L}}{\partial \mathbf{h}}
   & =\sum_{k=1}^{K}\frac{\partial \mathcal{L}}{\partial z_k}\frac{\partial z_k}{\partial \mathbf{h}}
  =\sum_{k=1}^{K}(p_k-\mathbb{1}[k=y])\,\mathbf{w}_k
  \label{eq:dh_ja}
\end{align}
が得られる。特に \eqref{eq:dh_ja} は「ロジット誤差がクラス重み $\mathbf{w}_k$ を介して再合成され、表現 $\mathbf{h}$ がどの方向へ更新されるか」を決定するため，
何が先に表現として学習されるかは主に \eqref{eq:dh_ja} の構造で説明できる。

$K=68$ クラスがドメイン $d\in\mathcal{D}$ とアミノ酸 $a\in\mathcal{A}_d$ の組で表現できるとする：
\begin{equation}
  k \leftrightarrow (d,a),\qquad
  \mathcal{D}=\{\mathrm{arc},\mathrm{bac},\mathrm{euk}\},\qquad
  K=\sum_{d\in\mathcal{D}}|\mathcal{A}_d|=68.
\end{equation}
このときロジットと確率を $z_{d,a}$、$p_{d,a}$ と書く。
ドメインに属する確率の総和を
\begin{equation}
  P(d)=\sum_{a\in\mathcal{A}_d} p_{d,a}
  \label{eq:Pd_ja}
\end{equation}
と定義する。

ドメイン単位の集約logitsとして
\begin{equation}
  s_d=\log\sum_{a\in\mathcal{A}_d} e^{z_{d,a}}
  \label{eq:sd_ja}
\end{equation}
を導入する。すると分母は
\begin{equation}
  \sum_{k=1}^{K}e^{z_k}
  =\sum_{d\in\mathcal{D}}\sum_{a\in\mathcal{A}_d}e^{z_{d,a}}
  =\sum_{d\in\mathcal{D}} e^{s_d}
\end{equation}
と表せ、$\{s_d\}$ 上のsoftmaxは周辺確率 $P(d)$ に一致する：
\begin{equation}
  \frac{e^{s_d}}{\sum_{d'}e^{s_{d'}}}
  =
  \frac{\sum_{a}e^{z_{d,a}}}{\sum_{d'}\sum_{a'}e^{z_{d',a'}}}
  =
  \sum_{a}p_{d,a}
  =
  P(d).
  \label{eq:softmax_sd_ja}
\end{equation}
正解クラスが $y\leftrightarrow(d^\star,a^\star)$ に対応するとき（正解ドメイン $d^\star$），
\eqref{eq:ce_ja} と \eqref{eq:sd_ja} を用いた微分により、ドメイン集約ロジットに関する勾配は
\begin{equation}
  \frac{\partial \mathcal{L}}{\partial s_d}
  =
  P(d)-\mathbb{1}[d=d^\star]
  \label{eq:ds_ja}
\end{equation}
となる。学習初期において $P(d)\approx 1/|\mathcal{D}|=1/3$ とみなせば，
$\partial \mathcal{L}/\partial s_{d^\star}\approx -2/3$
$\partial \mathcal{L}/\partial s_{d\neq d^\star}\approx 1/3$ であり，
ドメインレベルには次元の低い（$|\mathcal{D}|=3$）かつ大きさが $O(1)$ の誤差信号が存在する。

次に，クラス重みが同一ドメイン内で共通成分を獲得しやすいという仮定を明示するため，
各クラス重みを
\begin{equation}
  \mathbf{w}_{d,a} = \mathbf{u}_d + \mathbf{v}_{d,a}
  \label{eq:decomp_ja}
\end{equation}
と分解する。ここで $\mathbf{u}_d$ はドメイン $d$ に共通な方向（共有成分），
$\mathbf{v}_{d,a}$ はアミノ酸ごとの残差（差分成分）である。
\eqref{eq:decomp_ja} を \eqref{eq:dh_ja} に代入すると，
\begin{align}
  \frac{\partial \mathcal{L}}{\partial \mathbf{h}}
   & =
  \sum_{d\in\mathcal{D}}\sum_{a\in\mathcal{A}_d}
  \bigl(p_{d,a}-\mathbb{1}[d=d^\star,a=a^\star]\bigr)\,(\mathbf{u}_d+\mathbf{v}_{d,a})
  \notag \\
   & =
  \underbrace{\sum_{d\in\mathcal{D}}
    \Bigl(\sum_{a\in\mathcal{A}_d}p_{d,a}-\mathbb{1}[d=d^\star]\Bigr)\mathbf{u}_d}_{\text{ドメイン成分}}
  \;+\;
  \underbrace{\sum_{d\in\mathcal{D}}\sum_{a\in\mathcal{A}_d}
    \bigl(p_{d,a}-\mathbb{1}[d=d^\star,a=a^\star]\bigr)\mathbf{v}_{d,a}}_{\text{アミノ酸差分成分}}
  \label{eq:dh_split_ja}
\end{align}
を得る。\eqref{eq:Pd_ja} を用いるとドメイン成分は
\begin{equation}
  \sum_{d\in\mathcal{D}} \bigl(P(d)-\mathbb{1}[d=d^\star]\bigr)\,\mathbf{u}_d
  \label{eq:dh_domain_ja}
\end{equation}
と簡約される。すなわち，表現 $\mathbf{h}$ に流れる更新信号は
（i）3つの方向 $\{\mathbf{u}_d\}$ とドメイン誤差 \eqref{eq:ds_ja} で記述される低次元成分と，
（ii）多数の残差方向 $\{\mathbf{v}_{d,a}\}$ を必要とする高次元成分に分かれる。
同一ドメインの $\mathbf{w}_{d,a}$ が似た方向を向く（すなわち $\mathbf{u}_d$ が支配的）場合，
\eqref{eq:dh_domain_ja} の寄与は同一ドメインのサンプル間で整列しやすく，
勾配が建設的に加算されて最適化が速く進む。

一方，アミノ酸差分成分（\eqref{eq:dh_split_ja} 第2項）は、同一ドメイン内で
$|\mathcal{A}_d|\approx 22\text{--}23$ 個の異なる方向 $\mathbf{v}_{d,a}$ を学習し，
それらを互いに区別できるように表現 $\mathbf{h}$ を整形する必要がある。
異なるアミノ酸 $a\neq a'$ は更新方向が異なるため，
サンプルごとに $\mathbf{v}_{d,a}$ 方向の勾配が競合（干渉）しやすい。
このため，アミノ酸の細かな分離はドメイン分離に比べて高次元で難しく、学習がドメインに比べて遅れるのではないかと考えられれる。

まとめると
68クラスの損失 \eqref{eq:ce_ja} は
表現勾配 \eqref{eq:dh_ja} はクラス重みを介して再合成されるため，
同一ドメイン内の重みに共有成分 $\mathbf{u}_d$ が形成されると
低次元かつ大きいドメイン誤差\eqref{eq:ds_ja}が \eqref{eq:dh_domain_ja} の形で強く現れる。
これに対し，アミノ酸識別は多数の残差方向 $\mathbf{v}_{d,a}$ を必要とし干渉が生じやすい。
以上より、ドメインの分離がアミノ酸の分離より先に学習されることと考える。

\paragraph{(2)生物学的観点}

\ref{sec:rinalmo_inst}で述べた通り、CLSは1サンプルの情報を集約するトークンである。

そのため、CLSの学習推移を解釈するには、CLSが集約する各塩基トークン表現の学習推移にも注目する必要がある。
各塩基トークンの推移を観察すると、表現の分離は概ね (1)塩基(2)塩基を軸としたstructure(3)structureを軸としたドメイン の順に進行している。
一方で、アミノ酸特異性を中心とした分離が層として明確に現れる傾向は確認できなかった。
塩基についてはトークンIDとして明確に区別されるため、低層で分離が生じやすい。
また構造に関しては、RoPEにより相対位置に依存した相互作用を表現しやすく、二次構造に対応する特徴が中層で立ち上がりやすい。

ここで重要なのは、なぜ構造特徴がドメイン分離の主要な手がかりとなり得たかである。

先行研究\cite{marck2002trnomics}でも示されるように、tRNA配列のドメイン識別にはドメイン特有のモチーフ分布や構造的特徴が重要であることが知られている。

従って学習が進むにつれ、本モデルはtRNA配列の構造的特徴を捉え、その結果としてドメイン差を自然に抽出した可能性がある。

以上の2つの観点からモデルの学習推移に関して一定の考察を行った。


\subsection{AttnLRP解析}
\subsubsection{層ごとのトークンタイプ別寄与率}

\paragraph{初期層}

まず、層の浅い段階（0〜10層付近）では、Relevanceの絶対量・比率のいずれにおいても各塩基トークンが大部分を占めていることが確認された。CLS の寄与は相対的に小さく、EOS に関しては全層を通じてほぼ無視できる水準であった。

この結果は、初期層においてモデルが主に各塩基トークンに局所的に付随する情報を処理している段階であることを示唆している。
すなわち、初期層ではまだ情報が統合されておらず、配列中の個々の位置情報が分散的に扱われている状態にあると解釈できる。

\paragraph{中間層から後半層}

層が進むにつれて、CLS トークンの Relevance は単調に増加し、とくに20層以降ではその増加が顕著となった。
絶対値の観点では、最終層に近づくにつれて CLS の Relevance が各塩基トークンのrelevanceの合計を上回り、比率においても CLS が主要な寄与源となることが示された。

この傾向は、Transformer における表現学習の一般的理解と整合的である。
中間層以降では，各塩基トークンに分散していた情報が段階的に CLSに集約され、系列全体を表す表現として統合されていく可能性がある。
本分類モデルでは最終的な分類判断を CLS表現のみに基づいて行うことから、AttnLRPによる解析はこの解釈を定量的に示唆している。
そのためモデル内部で
各塩基トークン：局所的・構造的特徴の保持

CLS：それらを統合した判断表現という機能分化が段階的に形成されていることを示唆する。

\paragraph{EOSの寄与特性}

EOS トークンは全層を通じて Relevance が小さく、予測への寄与は相対的に限定的であった。
従ってRiNALMo における EOS は系列終端を示すための制御的トークンとして機能しており，意味的・識別的情報の集約には関与しにくい可能性がある。
この点は、タスクに有用な情報がCLSに集約されやすいというTransformerの性質と整合的である。

\paragraph{Relevanceの層ごとの減少}
本解析において、層を遡るにつれて（出力層から入力層へ向かうにつれて）、relevanceの総和が減少する傾向が見られた。そのためrelevance総和の減少についてRiNALMoのモデル構造とAttnLRPの算出特性の観点から考察する。

Relevance減少の主要な原因は、Relevance算出に用いたInput $\times$ Gradient法の特性と、RiNALMoのTransformer BlockにおけるLayerNormの配置の相互作用にあると考えられる。

AttnLRPの実装では、Relevanceの算出に Input $\times$ Gradient 法（$R(x) = x \odot \frac{\partial y}{\partial x}$）を採用している。
線形変換 $y = Wx + b$ に対するこの手法の特性として、重み行列 $W$ に起因する寄与は入力 $x$ に分配される一方で、バイアス項 $b$ に起因する寄与（$b \odot \frac{\partial y}{\partial y}$）は入力 $x$ には分配されず、消失するという性質がある。

RiNALMoの \texttt{TransformerBlock} の実装を確認すると、以下のコードに示すように、残差結合（Residual Connection）の経路自体が LayerNorm を通過する構造となっている。

\par
\begin{lstlisting}[language=Python, caption=RiNALMo TransformerBlockの実装抜粋]
def forward(self, x, key_padding_mask=None, need_attn_weights=None):
    x = self.attn_layer_norm(x)  
    if self.use_flash_attn:
        mh_out, attn = self.mh_attn(x, ...)
    else:
        mh_out, attn = self.mh_attn(x, attn_mask=None, ...)
    x = x + self.residual_dropout_1(mh_out)  

    residual = x
    x = self.out_layer_norm(x)
    x = residual + self.residual_dropout_2(self.transition(x))
    return x, attn
\end{lstlisting}


Pre-Norm構成（例: LLaMA(\cite{jain2019attention})等）では、残差経路は正規化層を回避して加算される（$x_{l+1} = x_l + \text{Attn}(\text{Norm}(x_l))$）ため、情報はバイアス項の影響を受けずに伝播する経路を持つ。
しかし、RiNALMoの構成では、情報の幹となる経路が各層で必ず LayerNorm を通過する。
LayerNorm には学習可能なバイアスパラメータ（$\beta$）が含まれており、層を重ねるごとにRelevanceの一部がこのバイアス項へと吸収され続けることになる。

RiNALMoは33層という深い層構造を持つため、この各層でのわずかなRelevanceのバイアスへの吸収が蓄積し、結果として入力層におけるRelevanceの総和が、出力層のRelevanceと比較して顕著に減少したと考えられる。
そのため、relevanceの値について考察する場合はその点に留意する必要があると考えられる。

\subsubsection{層ごとのコンポーネント割合}

本節では、AttnLRP により算出した Relevance を attention / FFN / residual の各コンポーネントに分解し、
層方向における寄与構造の変化を分析する。

\paragraph{CLSトークンにおける寄与構造の遷移}

\begin{itemize}
  \item \textbf{初期層（layer0--2）}では、AttentionおよびFFN成分が大きな正の寄与を示す一方で、
        Residual成分は負の寄与を持つ。
        この負のResidualは、残差経路が寄与していないことを意味するのではなく、
        CLSに蓄積された表現の一部が、最終予測に対して抑制的（反証的）に作用し、
        その影響が他成分（Attention/FFN）の正の寄与によって相殺されている可能性を示す。
        したがってこの層域は、CLSが単に前層表現を保持する段階というより、
        Attentionによる入力トークンからの情報流入と、FFNによる非線形変換を通じて
        CLSの初期表現を立ち上げる過程にあり、その際に残差経路に含まれる成分の取捨選択
        （不要成分の抑制）が強く生じている段階と解釈できる。

  \item \textbf{初期層後半（layer3--8）}では、Attention成分の比率が低下し、
        Residual成分が支配的となる。
        これは、初期層で形成されたCLS表現が以降の層で維持・蓄積される局面へ移行し、
        層をまたいだ表現の安定化が進むことを示唆する。
        また内部表現解析においても，この層域からドメインごとの分離（クラスタ形成）が
        徐々に観測されたことからResidual経路による表現保持が、
        ドメイン識別に有用な表現の固定化と対応している可能性がある。

  \item \textbf{中間層（layer8--19）}では、Residual成分が全体の約60--70\%を占め，
        AttentionおよびFFNの寄与は相対的に小さい。
        このことは、CLS表現が主として残差経路を通じて保持されながら、
        AttentionとFFNによって緩やかな更新（微修正・再符号化）を受ける段階にあることを示す。
        すなわち，この層域では大きく作り直すよりも保持しつつ整える更新様式が優勢であり，
        ドメイン情報が安定的に埋め込まれていると解釈できる。

  \item \textbf{後半層（layer20--33）}では、Residual成分の優位は維持される一方で、
        Attention成分の比率が再び増加する。
        内部表現解析ではこの層域で表現の主軸がドメインからアミノ酸へ移行していたことから，
        上位層ではCLSが注意機構を再度強めることでアミノ酸判別に必要なトークン情報を
        追加的に集約・選別している可能性が示唆される。
        ただし絶対値の増加が主としてResidualで駆動されている点は重要であり、
        最終判断に効くCLS表現が、各層で新規に生成されるというよりも、
        既存の表現が残差経路を通じて蓄積・増幅されつつ，
        上位層のAttentionがアミノ酸の分類に重要なのでないかと考えられる。
\end{itemize}




\paragraph{各塩基トークンにおける寄与構造の遷移}

\begin{itemize}
  \item \textbf{初期層（layer0--2）}では、layer0においてResidual成分がほぼ全体を占める一方、
        layer1--2ではAttentionおよびFFN成分が大きな正の寄与を示しResidual成分は負の寄与を持つ。
        これはCLSと同じ挙動を示しており
        各塩基トークン表現が埋め込み直後（layer0）には主に残差として伝播しているのに対し、
        layer1--2では、元の塩基トークン表現（Residual経路に含まれる成分）の一部が最終予測に対して抑制的に作用し、
        その影響をAttention/FFNによる更新が上書き・相殺しながら、
        表現を立ち上げている段階であることを示唆する。

  \item \textbf{初期層後半（layer3--5）}では、Residual成分が正に転じて急速に増加し、
        比率としてもResidualが過半を占めるようになる。
        これは，初期層で構築された文脈化表現が，残差経路を通じて安定的に保持・蓄積され始める
        フェーズへの移行を表すと解釈できる。

  \item \textbf{中間層（layer6--23）}では、Absoluteの総量が層とともに増加し、
        おおむねlayer21--23付近でピークを迎える。
        比率としてはResidualが概ね\(\sim 55\) -- \(75\%\)程度を占め，
        FFNが\(\sim 15\) -- \(30\%\)、Attentionが\(\sim 10\) -- \(20\%\)で推移する。
        このことは、各塩基トークン表現が主としてResidual経路で維持されつつ、
        FFNによる非線形な再符号化とAttentionによるトークン間混合を通じて、
        タスクに有用な特徴が徐々に洗練されていることを示唆する。
        また総Relevanceが中間層で最大化する点はこの層からCLSへのrelevanceが急激に高まることを意味する。

  \item \textbf{後半層（layer24--33）}では、各塩基トークン絶対値総量が明確に減少へ転じ
        最終層に向けて低下する。
        比率としてはResidual優勢が維持される一方で、終盤（特に最終層）ではAttention比率がやや増加する。
        この挙動は、上位層において意思決定に直接効くRelevanceが各塩基トークン側から離れ、
        CLSへと集約されていく過程と整合的である。
        その上で最終層付近のAttention比率の上昇は、
        残された少数の重要トークンに対して重みづけを行いながら，
        最終的な根拠の選別・調整を行っている可能性を示唆する
\end{itemize}



\subsection{内部表現解析とAttnLRP解析を統合した全体考察}

内部表現解析および AttnLRP 解析を統合すると、
RiNALMo における表現学習は層方向に沿って一貫した階層構造を持つことが明らかとなった。

浅層では、CLS・各塩基トークンともにラベル非依存な汎用表現を形成しており、
AttnLRP においても Relevance は主に各塩基トークンに分配されている。
この段階では、配列中の局所的・位置依存的特徴が分散的に処理されていると解釈できる。

中層に進むにつれて、内部表現解析では
CLS 表現においてドメインごとの分離が顕在化し、
各塩基トークン表現では structure を軸としたクラスタ形成が進行する。
AttnLRP 解析では、この層域において CLS への Relevance 集約が進み、
residual 成分を中心として情報が安定的に保持されていることが示された。
これは、ドメインおよび構造情報が中間概念として内部表現に組み込まれる段階に対応する。

さらに深層では、CLS 表現の主軸がドメインからアミノ酸へと移行し、
ドメインを横断した同一アミノ酸間の類似性が高まる。
一方、各塩基トークン表現は structure 内の細粒度な部位差へと分解される。
AttnLRP においても、この段階で CLS の attention 成分が再び増加しており、
各塩基トークンからアミノ酸識別に必要な情報を再集約していることが示唆される。

最終層では、CLS 表現がタスクに特化した判断表現として機能し、
分類判断に必要な情報が集中的に集約される。
このことは、Relevance が最終的に CLS に強く集中する挙動とも整合的である。

以上より、RiNALMo は
\begin{enumerate}
  \item 浅層で局所的特徴を分散的に抽出し、
  \item 中層で structure・ドメインという中間表現を形成し、
  \item 深層でドメインを超えたアミノ酸中心表現へ再編成した上で、
  \item CLS によって最終判断を行う
\end{enumerate}
という階層的かつ再帰的な表現学習機構を持つことが示された。
内部表現解析と AttnLRP 解析は、この過程をそれぞれ一貫して支持している。


\subsection{単一サンプルにおけるrelevance解析}

本研究では全ラベルを対象として解析を実施すると結果が膨大となるため、bac\_Ala、arc\_Ala、euk\_Alaの代表サンプルに限定し、層ごとのトークン別relevanceの可視化結果およびrelevanceに対する検定結果を示す。

\ref{fig:attnlrp_arc_Ala}、\ref{fig:attnlrp_bac_Ala}、\ref{fig:attnlrp_euk_Ala}に、それぞれarc\_Ala、bac\_Ala、euk\_Alaの層ごとのトークン別relevanceを示す。なお、各図では累積relevanceが50\%に達するまでのトークンのみを表示している。

フローグラフより、いずれのサンプルおよび層においても、少数の特定塩基がrelevanceの大部分を占めるという明確な傾向は観察されなかった。
すなわち、relevanceは特定箇所に強く集中するというよりも、配列全体に比較的広く分散していることが示唆される。
一方で、いずれのサンプルにおいてもアンチコドン領域（図の中央付近）は全層にわたり相対的に高いrelevanceを示す傾向が認められた。
これはモデル構造上、CLSがアンチコドン領域へ注意を向けやすく、層を遡っても当該領域のrelevanceが相対的に高く維持される可能性によるものと考えられる。

以上の結果から，本モデルは分類に際して特定の塩基に過度に依存するのではなく，配列全体に分散した情報を統合して利用している可能性が示唆される。
先行研究では，ドメインやアミノ酸残基ごとに決定因子が存在すると報告されている一方で，本モデルでは局所的特徴への依存が主要因とはならず，広域的な文脈情報を活用していると解釈できる。

ただし、「特定の決定因子が顕在化しない」という結果は、本モデルのタスク設定に依存する可能性がある。

本モデルは \texttt{domain\_amino\_acid} を分類ラベルとしており、特定の結合モチーフ検出のような局所特徴中心のタスクとは異なる。
以上の解析から分かる通り、モデルはドメイン、アミノ酸をそれぞれ段階を踏んで分類していると考えられ、
特にドメイン分類においては、配列全体を統合的に取り扱う必要がある可能性が高い。
従って、決定因子が明確に現れにくいことは解析の不備ではなく、モデルの性質の問題の可能性も挙げられる。

さらに、AttnLRPの性質を踏まえた解釈が必要である。AttnLRP は，モデル出力（分類ヘッドの logit）を起点として信号を逆伝播させ、各塩基トークンのrelevanceを推定する手法である。

\begin{enumerate}
  \item \textbf{出力から \(\mathrm{CLS}\) への逆伝播:}
        分類ヘッドの logit を起点として逆伝播を行い，最終層の \(\mathrm{CLS}\) 塩基トークンに割り当てられる relevance を算出する。

  \item \textbf{層方向の逆伝播:}
        続いて Transformer の各層を上位層から順に遡って逆伝播する。理想的には，層 \(N\) における relevance \(R_N\) と層 \(N-1\) における relevance \(R_{N-1}\) は（近似的に）保存され，\(R_N\) は \(R_{N-1}\) から構成されるとみなせる。
\end{enumerate}

ただし，Transformer を構成する各コンポーネント（Attention,FFN,Residual）の性質を踏まえると、relevance の解釈において次の点が重要となる。
\begin{itemize}
  \item \textbf{FFN:}
        非線形変換を通じて特徴を変換・統合するが、内部でどの特徴がどのように合成されているかは観測しにくい。そのためFFN による変換を経た relevance は，入力特徴へ単純に対応づけられない場合がある。

  \item \textbf{Attention:}
        他トークンから情報を収集・統合する機構である。ここで重要なのは，層 \(N\) で高い relevance を示すトークンが，必ずしも層 \(N-1\) でも高い relevance を示すとは限らない点である。すなわち，層 \(N\) で重要とされる表現が，層 \(N-1\) の別トークンから Attention により集約された結果である可能性がある。逆に，層 \(N\) で重要なトークンが，下位層へ遡る過程で relevance を他トークンへ分配している可能性もある。

  \item \textbf{Residual（残差接続）:}
        表現が層をまたいで加算的に伝搬するため、relevance も複数経路に分散し得る。結果として，「どの層・どの経路を通って重要性が形成されたか」を一意に特定しにくい。
\end{itemize}

以上より，入力層まで逆伝播した際に高い relevance を示すトークンを，ただちに「決定因子」と結論づけることはできない。層を遡る過程で Attention や FFN によって情報が再分配・再統合され得るためである。

また，近年の言語モデル解釈は，トークン単体の重要度を推定する枠組みが主流であり，二次・三次の相互作用(トークンの相互効果)までを逆伝播により安定に推定する一般的手法は限定的である。

したがって，本解析のみから入力側の高 relevance 塩基を決定因子として断定することは避けるべきである。

ただ本解析において特に重要な観察として、23層付近からモデル挙動が変化する点が挙げられる。既述のとおり、本モデルでは23層付近からCLSへの寄与が急増する傾向が確認されている。この事実は、モデル内部において情報処理のフェーズが切り替わっている可能性を示唆する。

\begin{itemize}
  \item \textbf{23層から最終層:} 当該区間では、relevanceの高い塩基に関する情報がCLSへ集約される過程にある可能性が高い。従って、この区間で高いrelevanceを示す塩基は、最終判断により直接的に関与している可能性がある。
  \item \textbf{Inputから23層:} それ以前の層におけるrelevance分布は、最終的な集約に至る前段階としての情報移動や中間表現形成過程を反映している可能性がある。
\end{itemize}

以上より、全層を一様に評価するのではなく、(i) Inputから23層、(ii) 23層から最終層、の二区間に分けてrelevanceを比較することが有効であると考えられる。特に、23層以降でrelevanceが増大する、または高値を維持する塩基がモデル判断にとって相対的に重要である、という仮説の下で結果を解釈・提示することが妥当である。

表\ref{tab:token_arc_stats}、\ref{tab:token_bac_stats}、\ref{tab:token_euk_stats}に、それぞれarc\_Ala、bac\_Ala、euk\_Alaの各層におけるトークン別検定結果を示す。

arc\_Alaおよびbac\_Alaでは、アンチコドンに対応するC36が有意に高いrelevanceを示した。
一方、euk\_AlaではC36は有意ではなかった。
また、本モデルの解釈に基づけば、アンチコドンとして不変な塩基であるC35も有意となることが予想されたが、いずれのサンプルにおいても有意差は確認されなかった。

さらに、arc\_Alaおよびbac\_Alaでは、Alaの重要な決定因子として知られるT70が有意に高いrelevanceを示した。
一方で、euk\_Alaでは同様の傾向は確認されなかった。
これは、本モデルがT70のような単一の既知決定因子に依存せず、配列全体に分散した特徴を用いて分類している可能性を示唆する。

領域ごとの傾向として、arc\_Alaではvariable regionに対して有意に高いrelevanceを示す傾向が認められた。
bac\_AlaではD armおよびacceptor armにおいて有意に高いrelevanceを示す傾向が観察された。
euk\_Alaではアンチコドン領域で高いrelevanceを示す傾向が認められた。

以上を総括すると、本モデルにより得られた重要塩基は既知の決定因子と一致する部分も存在する一方で、多くの塩基については既知の決定因子と一致しない結果となった。

これに関しては
\ref{sec:transision_analysis}で考察したように本モデルには学習推移が存在することや
\ref{Consideration-top10-misclassifications}で考察したようにaaRS,tRNAの水平伝播の影響を考慮すると先行研究で示されているような決定因子が必ずしも本モデルの判断根拠として顕在化しないのではないかと考えられる。

これらのことから、他の解釈方法を用いた決定因子の推定や、解析手法自体の発展も望まれる。



\clearpage
\section{結論}

本研究では、RNA言語モデルRiNALMo\cite{penic2025rinalmo}を用いたtRNAの分類モデルを構築し、その内部表現と判断根拠の解明を試みた。
構築したモデルは、3ドメインにわたるデータセットに対する評価（テストデータ）において、Accuracy 0.9294、Weighted-F1 0.9286という極めて高い分類性能を達成した。

内部表現解析およびAttnLRPを用いた寄与度解析により、これまでブラックボックスとされてきた深層学習モデルの判断プロセスについて、以下の知見が得られた。

第一に、層ごとの学習推移が明確に確認された。
CLSは、中層（20層付近）でのドメイン情報の獲得を経て、深層（23層以降）では生物種を超えたアミノ酸特異的な表現へと段階的に遷移することが確認された。
一方、各塩基トークンの表現は、主にtRNAの二次構造を反映して分離し、深層ではより微細な構造単位へと分化する傾向を示した。
これにより本モデルが階層的な学習を自立的に行っていることが確認された。
またこの学習推移に関してはモデルの構造と生物学的観点から解釈を行っている。

第二に、AttnLRPによる解析から、情報の集約プロセスが定量的に示された。
初期層では各配列トークンに分散していたRelevanceが、層が進むにつれてCLSトークンへと集約される挙動が確認された。
また、特定の決定因子が高い寄与を示す一方で、判断根拠は特定塩基のみに依存せず配列全体に分散していることが示唆された。
これは本モデルの特性や、水平伝播の影響によるものだと考えられる。

一方で、課題として
(1)少数クラスに対する精度の低下(2)AttnLRPにおけるRelevance減衰の影響などが挙げられる。

そのため今後は、データ拡張等による不均衡データの克服や、AttnLRP以外の解釈手法を用いた解析や解析手法自体の改良を通じて、より精緻なモデル解釈と性能向上を目指すことが重要であると考えられる。

\clearpage
\section{付録}

\subsection{ソースコード、python環境}
\label{subsec:source-code}

python環境に関してはそれぞれのツールで推奨されているconda環境を用いている。

\begin{table}[H]
  \centering

  \label{tab:tools}
  \begin{tabular}{ll}
    \hline
    Tool        & URL                                                  \\
    \hline
    RiNALMo     & \texttt{https://github.com/lbcb-sci/RiNALMo}         \\
    cd-hit      & \texttt{https://github.com/weizhongli/cdhit}         \\
    tRNAscan-se & \texttt{https://github.com/UCSC-LoweLab/tRNAscan-SE} \\
    \hline
  \end{tabular}
  \caption{本研究で用いたpython環境のgithubリポジトリ}
\end{table}


\begin{table}[H]
  \centering
  \small
  \setlength{\tabcolsep}{5pt}
  \begin{tabularx}{\linewidth}{@{}>{\ttfamily}X X@{}}
    \toprule
    ファイル                                                   & 役割                                    \\
    \midrule
    Code/data/data\_pretreatment.py                        & データの前処理                               \\
    Code/data/cd-hit.py                                    & cd-hitの実行                             \\
    Code/data/tRNAscan-se.py                               & tRNAscan-seの実行                        \\
    Code/data/finetuning.py                                & ファインチューニング用のデータ分割                     \\
    Code/model/extract\_all\_layers\_embedding.py          & 分類モデルを用いたembedding抽出                  \\
    Code/model/model\_verification.py                      & 評価指標の計算と結果出力                          \\
    Code/model/predict\_label.py                           & 分類モデルを使った正答ラベルの抜き出し                   \\
    Code/Centered\_Cosine\_Similarity/CCS\_heatmap\_CLS.py & CLS の コサイン類似度 ヒートマップの作成               \\
    Code/UMAP/umap\_CLS.py                                 & CLS の UMAP 可視化                        \\
    Code/UMAP/umap\_token.py                               & 各塩基トークン の UMAP 可視化                    \\
    AttnLRP/lxt\_rinalmo\_patches.py                       & RiNALMoに対するAttnLRPの適用                 \\
    AttnLRP/add\_relevance\_features.py                    & AttnLRP を用いた relevance 計算・保存          \\
    AttnLRP/analyze\_layer\_contribution.py                & 層ごとの塩基トークンタイプの可視化                     \\
    AttnLRP/analyze\_component\_breakdown.py               & 層ごとのコンポーネントの可視化                       \\
    AttnLRP/generate\_component\_flow\_html.py             & 単一配列における relevance の流れ可視化、relevance検定 \\
    rinalmo/model/attention.py                             & AttnLRPのため内部構造変更(モデル本体)               \\
    rinalmo/model/modules.py                               & AttnLRPのため内部構造変更(モデル本体)               \\
    rinalmo/model/model.py                                 & AttnLRPのため内部構造変更(モデル本体)               \\

    \bottomrule
  \end{tabularx}
  \caption{ソースコードの所在と役割（リポジトリ: \repo）}
\end{table}

またRiNALMoのファインチューニングのコードに関しても\texttt{https://github.com/lbcb-sci/RiNALMo}を参照してほしい。

\subsection{モデルパラメータ}
\label{subsec:model-params}

\begin{table}[H]
  \centering
  \begin{tabular}{ll}
    \hline
    項目                                            & 設定                                          \\
    \hline
    Backbone LM                                   & RiNALMo (\texttt{lm\_config}=\texttt{giga}) \\
    Transformer embedding dim ($c_{\mathrm{in}}$) & 1280                                        \\
    Head hidden dim ($d_{\mathrm{head}}$)         & 256                                         \\
    \#classes ($C$)                               & 68                                          \\
    Loss                                          & Cross-entropy                               \\
    Learning rate                                 & $8\times 10^{-6}$                           \\
    Weight decay                                  & $1\times 10^{-2}$                           \\
    \hline
  \end{tabular}
  \caption{RiNALMoのファインチューニングパラメータ}
\end{table}

\begin{table}[H]
  \centering
  \begin{tabular}{ll}
    \hline
    Parameter              & 設定                     \\
    \hline
    $n_{\text{neighbors}}$ & 20(CLS) or 30(各塩基トークン) \\
    metric                 & cosine                 \\
    min\_dist              & 0.1                    \\
    \hline
  \end{tabular}
  \caption{UMAPのパラメータ}
\end{table}



\subsection{データセット}
\label{subsec:datasets}
\begin{table}[H]
  \centering
  \begin{tabular}{ll}
    \hline
    ファイル                              & 役割                       \\
    \hline
    data/formatting/final.csv         & ファインチューニングに用いるデータセット     \\
    data/formatting/pretreatment.csv  & 基本的な前処理後のデータセット          \\
    data/analysis/origin/analysis.csv & 内部表現解釈、AttnLRPで用いるデータセット \\
    \hline
  \end{tabular}
  \caption{データセットの所在と役割（リポジトリ: \repo）}
\end{table}

\subsection{参考文献}
\bibliographystyle{junsrt}
\bibliography{refs}

\clearpage
\section{謝辞}
本研究及び本論文の執筆にあたり、ご指導をしてくださった北海道大学大学院 情報
科学研究科 生命人間情報科学部門 バイオインフォマティクス分野 情報生物学研究
室、遠藤俊徳教授ならびに長田直樹准教授、そして共に研究に励み、研究活動を
支えていただいた情報生物学研究室の皆様、平山陽子様に深く感謝いたします。


\end{document}





